% ======================================================================
% TopoSHAP -- ICLR 2027 submission (skeleton)
%
% Build:  make          (pdflatex + bibtex, output in build/)
% Clean:  make clean
%
% NOTE: iclr2027_conference.sty is a PLACEHOLDER; swap for the official
% ICLR 2027 style file (and its .bst) when released.
% ======================================================================
\documentclass[10pt,letterpaper]{article}

\usepackage{iclr2027_conference}   % PLACEHOLDER style -- see file header
% ======================================================================
% preamble.tex -- packages, theorem environments, draft-mode machinery
% ======================================================================

% ---------------------------------------------------------------- fonts
% (guarded so the skeleton also compiles under XeTeX/LuaTeX engines,
% e.g. tectonic; pdflatex behavior is unchanged)
\ifdefined\XeTeXrevision\else\ifdefined\directlua\else
  \usepackage[utf8]{inputenc}
  \usepackage[T1]{fontenc}
\fi\fi
\IfFileExists{microtype.sty}{\usepackage{microtype}}{}

% ------------------------------------------------------------- packages
\usepackage{booktabs}
\usepackage{graphicx}
\usepackage{amsmath}
\usepackage{amssymb}
\usepackage{amsthm}
\usepackage{xcolor}
\usepackage{colortbl}
\usepackage{xspace}
\usepackage{natbib}          % official ICLR style ships its own .bst; see main.tex

% Boxed-theory environment: tcolorbox if available, framed as fallback.
\IfFileExists{tcolorbox.sty}{%
  \usepackage{tcolorbox}%
  \newtcolorbox{theorybox}{%
    colback=gray!6, colframe=black!70, boxrule=0.6pt, arc=1.5mm,
    left=7pt, right=7pt, top=6pt, bottom=6pt}%
}{%
  \usepackage{framed}%
  \newenvironment{theorybox}{\begin{framed}}{\end{framed}}%
}

% hyperref + cleveref last
\usepackage[colorlinks=true,
            linkcolor=blue!55!black,
            citecolor=green!40!black,
            urlcolor=blue!55!black]{hyperref}
\usepackage[capitalize,noabbrev]{cleveref}

% --------------------------------------------------- theorem environments
\theoremstyle{plain}
\newtheorem{theorem}{Theorem}
\newtheorem{proposition}{Proposition}
\newtheorem{lemma}{Lemma}
\newtheorem{corollary}{Corollary}
\theoremstyle{definition}
\newtheorem{definition}{Definition}
\newtheorem{assumption}{Assumption}
\theoremstyle{remark}
\newtheorem{remark}{Remark}

% Fixed-label statements for the boxed theory blocks (T-numbered, so they
% survive reordering of the ordinary theorem counter).
\newtheorem*{propositionT}{Proposition~T1}
\newtheorem*{corollaryT}{Corollary~T2}
\newtheorem*{propositionTrank}{Proposition~T3}

% ======================================================================
% Draft-mode machinery
% ======================================================================

% ----- frozen-results toggle ------------------------------------------
% While results are UNFROZEN (handoff-sourced numbers, not yet regenerated
% from data/frozen/), every number referenced via \unfrozen{...} carries a
% gray "UNFROZEN" superscript. Flip the boolean below once
% scripts/make_tables.py has regenerated results_macros.tex from
% data/frozen/ -- all markers disappear and numbers render plainly.
\newif\ifFrozenResults
\FrozenResultstrue   % results frozen 2026-08-18 (manifest in data/frozen/MANIFEST.json)

\newcommand{\unfrozen}[1]{%
  \ifFrozenResults
    {#1}%
  \else
    {#1}\textsuperscript{\,\textcolor{gray}{\tiny\textsf{UNFROZEN}}}%
  \fi}

% ----- researcher TODO margin notes -----------------------------------
% \todoresearcher{...}       margin note (body text only)
% \todoresearcherinline{...} inline note -- REQUIRED inside floats and
%                            captions, where \marginpar is illegal.
\newcommand{\todoresearcher}[1]{%
  \marginpar{\raggedright\scriptsize\textcolor{red!70!black}{%
    \textbf{TODO(researcher):} #1}}}
\newcommand{\todoresearcherinline}[1]{%
  \textcolor{red!70!black}{\scriptsize\textbf{TODO(researcher):} #1}}

% ----- misc helper macros ---------------------------------------------
\newcommand{\method}{TopoSHAP\xspace}
\DeclareMathOperator*{\argmax}{arg\,max}

% Bolded one-sentence claim opening an experiment subsection.
\newcommand{\expclaim}[1]{\noindent\textbf{#1}\par\smallskip}

% Boxed practitioner's recipe: plain-language numbered steps, no math.
\newenvironment{recipebox}{\begin{theorybox}\noindent
  \textbf{Practitioner's recipe.}\begin{enumerate}\setlength{\itemsep}{1pt}}
  {\end{enumerate}\end{theorybox}}

% Neighborhood names (message-passing routes); always "neighborhoods",
% never "relations".
\newcommand{\nbhd}[1]{\texttt{#1}}

% Placeholder box for figures whose PDF has not been generated yet.
\newcommand{\figureplaceholder}[1]{%
  \fbox{\parbox[c][0.30\linewidth][c]{0.95\linewidth}{%
    \centering\color{gray}\textsf{FIGURE PLACEHOLDER}\\[2pt]
    \texttt{\detokenize{#1}}}}}

% Convenience shorthand
\newcommand{\eg}{e.g.,\xspace}
\newcommand{\ie}{i.e.,\xspace}
\newcommand{\rk}{\operatorname{rk}}

% ======================================================================
% results_macros.tex -- ALL result numbers used anywhere in the paper.
%
% *** THIS FILE IS REGENERATED by scripts/make_tables.py from
% *** data/frozen/.  DO NOT EDIT NUMBERS BY HAND.
%
% *** CURRENT VALUES ARE HANDOFF-SOURCED AND UNFROZEN: they were copied
% *** from the experiment handoff notes and have NOT yet been
% *** regenerated from frozen data.  While this is the case, keep
% *** \FrozenResultsfalse in preamble.tex so every number rendered via
% *** \unfrozen{\ResFoo} carries the gray UNFROZEN marker.
%
% Convention: in section text, NEVER write a result number literally;
% always reference it as \unfrozen{\ResFoo}.
% ======================================================================

% ---------------------------------------------------------------- misc
% Visible stand-in for numbers that exist in the logs but were not part
% of the handoff; make_tables.py must replace every use of this macro.
\newcommand{\ResTBD}{\textcolor{red}{\textsc{tbd}}}

% ------------------------------------------------- Leg 1: predictions
\newcommand{\ResSurrogateAgreement}{0.9579}   % MANTRA: fraction of model predictions reproduced by counting surrogate
\newcommand{\ResHodgeAccuracy}{0.8764}        % MANTRA orientability accuracy with Hodge spectral features
\newcommand{\ResOrientStruct}{0.541}          % MANTRA orientability with structural-only features (control)
\newcommand{\ResCountingCeiling}{0.5054}      % MANTRA orientability balanced acc. of count-only models (record #6). NOT an information ceiling: oracle-thresholded counts reach 0.799 balanced, count-rank AUROC 0.784
\newcommand{\ResGEFBenzene}{\ResTBD}          % GEF faithfulness, Benzene            % TODO(researcher): fill from frozen data
\newcommand{\ResGEFFluoride}{\ResTBD}         % GEF faithfulness, Fluoride-Carbonyl  % TODO(researcher): fill from frozen data
\newcommand{\ResGEFMutagenicity}{\ResTBD}     % GEF faithfulness, Mutagenicity       % TODO(researcher): fill from frozen data
\newcommand{\ResGEFSubgraphXFluoride}{\ResTBD}% SubgraphX GEF, Fluoride-Carbonyl     % TODO(researcher): fill from frozen data
\newcommand{\ResGEFSubgraphXMutagenicity}{\ResTBD} % SubgraphX GEF, Mutagenicity     % TODO(researcher): fill from frozen data

% ---- Table 1 (attribution quality) values from the verified claims
% ---- ledger (data/frozen/reference/synthesis_report/claims_ledger.json).
% ---- GEA = node-projected ground-truth explanation agreement (higher
% ---- better); GEF differences are confidence-matched paired deltas
% ---- (ours minus flat GNN; negative = more faithful). UNFROZEN.
\newcommand{\ResGEAOursBenzene}{0.725}        % sgx_ours_Benzene
\newcommand{\ResGEAOursFluoride}{0.637}       % sgx_ours_FluorideCarbonyl
\newcommand{\ResGEAOursMutagenicity}{0.744}   % sgx_ours_Mutagenicity
\newcommand{\ResGEASubgraphXBenzene}{0.752}   % sgx_ref_Benzene (strongest SubgraphX reference)
\newcommand{\ResGEASubgraphXFluoride}{0.316}  % sgx_ref_FluorideCarbonyl
\newcommand{\ResGEASubgraphXMutagenicity}{0.394} % sgx_ref_Mutagenicity
\newcommand{\ResGEARandomBenzene}{0.238}      % benzene_chance (uniform-random control GEA)
\newcommand{\ResSGXAboveChanceWins}{7}        % SubgraphX beats the random control in ...
\newcommand{\ResSGXAboveChanceCells}{8}       % ... of these dataset-model combinations
\newcommand{\ResGEFDiffBenzeneGCN}{-0.878}    % gef_cm_Benzene_gcn
\newcommand{\ResGEFDiffBenzeneGIN}{-0.427}    % gef_cm_Benzene_gin
\newcommand{\ResGEFDiffMutagenicityGCN}{-0.079} % gef_cm_Mutagenicity_gcn
\newcommand{\ResGEFDiffMutagenicityGIN}{-0.227} % gef_cm_Mutagenicity_gin
\newcommand{\ResGEFDiffAlkaneGCN}{-0.094}     % gef_cm_AlkaneCarbonyl_gcn
\newcommand{\ResGEFDiffAlkaneGIN}{-0.483}     % gef_cm_AlkaneCarbonyl_gin

% ------------------------------------- Leg 2: architecture selection
\newcommand{\ResNumNeighborhoods}{11}         % candidate neighborhoods in the GCCN game
\newcommand{\ResStemEpochs}{5}                % epochs of all-neighborhood ensemble stem training
\newcommand{\ResSampledPasses}{128}           % forward passes for the sampled Shapley estimator
\newcommand{\ResRecipeCostMin}{1.2}           % guided-recipe cost, min (measured training-run equivalents)
\newcommand{\ResRecipeCostMax}{3.6}           % guided-recipe cost, max (measured training-run equivalents)
\newcommand{\ResSampledCostFraction}{$\sim$1/10} % sampled-game cost relative to exact game
\newcommand{\ResRegretMax}{0.024}             % max accuracy regret of sampled-game picks vs exact-game picks
\newcommand{\ResRhoExactCheap}{0.8322}        % Spearman rho, exact-game vs cheap (sampled) attributions
\newcommand{\ResGuidedWins}{34}               % guided-beats-random wins ...
\newcommand{\ResGuidedTrials}{42}             % ... out of trials
\newcommand{\ResWarmColdDelta}{+0.0024}       % warm-start (continue from stem) minus cold-retrain accuracy delta
\newcommand{\ResAccHeldWins}{6}               % datasets where the accuracy claim held ...
\newcommand{\ResAccHeldTotal}{7}              % ... out of datasets
\newcommand{\ResNStarFluoride}{1.2}           % N* on Fluoride-Carbonyl (stated exception)
\newcommand{\ResNStarNCI}{1.5}                % N* on NCI1 (stated exception)
\newcommand{\ResKDefault}{6}                  % safe default k
\newcommand{\ResKSmall}{3}                    % k winning on Mutagenicity / citation graphs
\newcommand{\ResSweepEqMin}{4.7}              % REFUTED exact-game pricing, min sweep-equivalents
\newcommand{\ResSweepEqMax}{28.1}             % REFUTED exact-game pricing, max sweep-equivalents
\newcommand{\ResPooledRho}{0.3212}            % pooled rho, Shapley reliance vs post-pruning viability
\newcommand{\ResPooledRhoCILow}{0.057}        % pooled rho 95% CI, lower
\newcommand{\ResPooledRhoCIHigh}{0.605}       % pooled rho 95% CI, upper
\newcommand{\ResMutagStability}{0.94}         % MUTAG cross-seed stability of Shapley rankings

% ----------------------------------------------- framework validation
\newcommand{\ResExactnessMaxErr}{$2.254\times 10^{-14}$} % max |exact - brute force| deviation

% -------------------------------------------------- trustworthiness
\newcommand{\ResNumBugsFramework}{3}          % TopoBench framework bugs caught by axiom checks
\newcommand{\ResNumBugsTotal}{4}              % including the composed cross-rank batching bug (invariance guard)

% ======================================================================
% Attribution-ladder pipeline (frozen rule 2026-08-17) and its evidence.
% Source: docs/method-improvements.md + the 2026-08-17 framework brief.
% ALL VALUES UNFROZEN (handoff/dev-log sourced) until make_tables.py
% regenerates them from data/frozen/.
% ======================================================================

% ------------------------------------------ pipeline costs (budget dial)
\newcommand{\ResLadderCostBOne}{1.3}          % B=1 single-shot cost, ~training-run equivalents
\newcommand{\ResLadderCostBTwo}{2.1}          % B=2 reliability tier cost, ~training-run equivalents
\newcommand{\ResMenuEntries}{6}               % literature menu = B=6 with literature rungs

% ------------------------- out-of-sample validation of the frozen rule
\newcommand{\ResLadderOOSContained}{4}        % OOS landscapes where the ladder contains the exact optimum ...
\newcommand{\ResLadderOOSTotal}{5}            % ... out of never-touched exact landscapes
\newcommand{\ResLadderOOSWorstRegret}{+0.0068}% worst OOS regret (ZINC-GIN landscape)
\newcommand{\ResLadderAllContained}{10}       % exact-zero-regret landscapes ...
\newcommand{\ResLadderAllTotal}{12}           % ... out of all exact landscapes
\newcommand{\ResLadderBOnePrice}{+0.022}      % worst B=1 (stop-rung-only) regret, ZINC: measured price of budget 1

% ---------------------------------------------- automatic k / stop rule
\newcommand{\ResBinomOrderDatasets}{7}        % datasets where binomial floor matches measured retraining noise in order of magnitude (all 7)
\newcommand{\ResBinomConservativeWins}{6}     % datasets where the floor errs conservative (slightly larger) ...
\newcommand{\ResBinomConservativeTotal}{7}    % ... out of marker-suite datasets
\newcommand{\ResGreedyStopBinomRegret}{0.0090}% mean regret, greedy-stop with 1-se binomial threshold, scored on exact truth
\newcommand{\ResGreedyStopOracleRegret}{0.0121}% mean regret, greedy-stop with oracle (measured-noise) threshold
\newcommand{\ResFixedKMeanRegret}{0.0313}     % fixed k=3 (top-k by phi) mean regret on the 12 exact landscapes
\newcommand{\ResFixedKMutagRegret}{0.2057}    % fixed k=3 max regret (MUTAG catastrophe)
\newcommand{\ResMutagTrueKStar}{5}            % MUTAG's true optimal coalition size k*

% ----------------------------- plateau finding / forward-greedy null
\newcommand{\ResPlateauZeroCells}{20}         % experiment cells with exactly-zero best forward gains ...
\newcommand{\ResPlateauTotalCells}{35}        % ... out of cells in the cross-cluster replication
\newcommand{\ResForwardFluoride}{0.848}       % forward-greedy test accuracy, Fluoride-Carbonyl (null)
\newcommand{\ResForwardFluorideRef}{0.936}    % reference arm test accuracy, Fluoride-Carbonyl
\newcommand{\ResForwardCoraMean}{0.769}       % forward-greedy test accuracy mean, Cora (null)
\newcommand{\ResForwardCoraSd}{0.081}         % forward-greedy test accuracy sd, Cora
\newcommand{\ResBackwardCoraSd}{0.008}        % backward-elimination test accuracy sd, Cora (variance collapse)
\newcommand{\ResBackwardBeatsForwardWins}{6}  % datasets where backward beats forward ...
\newcommand{\ResBackwardBeatsForwardTotal}{7} % ... out of marker-suite datasets

% ------------------- substitute-commitment replication (HOPSE family)
\newcommand{\ResHopseCommitDropPts}{5.3}      % accuracy points lost by the committed HOPSE pick (Mutagenicity, seed 43)
\newcommand{\ResHopseCommitValBad}{0.943}     % validation accuracy of the committed HOPSE pick ...
\newcommand{\ResHopseCommitValSibA}{0.992}    % ... vs sibling seeds' validation accuracies
\newcommand{\ResHopseCommitValSibB}{0.970}

% ---------------------------- catastrophe audit (validation visibility)
\newcommand{\ResCatastropheCount}{2}          % catastrophic picks across the development grid ...
\newcommand{\ResCatastropheCells}{63}         % ... out of grid cells
\newcommand{\ResCatastropheValA}{0.759}       % validation accuracy of catastrophic pick 1 (Benzene seed 43)
\newcommand{\ResCatastropheValB}{0.760}       % validation accuracy of catastrophic pick 2 (Benzene seed 43)
\newcommand{\ResCatastropheSiblingVal}{0.92}  % ~ validation accuracy of sibling picks
\newcommand{\ResCatastropheFalseAlarms}{0}    % false alarms at the same validation threshold

% ------------------------------------------- unique-eval cost accounting
\newcommand{\ResUniqueEvals}{129}             % ~ unique cached coalition evaluations behind the "128-pass" budget
\newcommand{\ResExactGameCoalitions}{2{,}047} % non-empty coalitions of the exact 11-player game
\newcommand{\ResUniqueEvalFraction}{0.06}     % ~ unique-eval count / exact-game count

% ------------------------------------------ baselines and the lead figure
\newcommand{\ResHopseSweepConfigs}{12}        % HOPSE paper-time sweep size (configs)
\newcommand{\ResGCCNSweepConfigs}{23}         % ~ GCCN neighborhood-sweep size (configs)
\newcommand{\ResTDLOriginalCount}{4}          % original TDL architectures shown as baselines
\newcommand{\ResTopoTuneRouteDefault}{4}      % TopoTune shipped default: routes in its config
\newcommand{\ResTopoTuneSweepStructures}{10}  % TopoTune paper-time sweep: neighborhood structures (x sub-models)
\newcommand{\ResHopseModelDefaultBlocks}{3}   % HOPSE shipped model default: encoding blocks
\newcommand{\ResHopseExptDefaultBlocks}{1}    % HOPSE shipped experiment default: encoding blocks
\newcommand{\ResBenzeneGCCNCeiling}{0.9252}   % Benzene: exact ceiling of the entire GCCN neighborhood space
\newcommand{\ResBenzeneSpectralMargin}{0.01}  % ~ margin by which HOPSE spectral encodings clear that ceiling
\newcommand{\ResMoleculeDatasets}{6}          % molecule datasets carrying the HOPSE-family markers (ladder + fixed k=3 rung)
\newcommand{\ResMarkerDatasets}{7}            % datasets in the (development) marker suite

% ---- Table 2 (performance) headline aggregates, computed from
% ---- data/frozen/{autok_marker,ladder_b2_{marker,t1,mantra},hopse_k3_rung,
% ---- baseline_markers,hopse_sweep_{t1,mantra}} by scripts/make_tables.py
% ---- --performance-table (which also emits the table body,
% ---- paper/tables/performance_summary.tex, and patches every macro in
% ---- this block plus the campaign-11 block below). UNFROZEN.
\newcommand{\ResSweepCostMultMin}{10}         % GCCN-sweep FLOPs / ours-B=1 FLOPs, min over 7 datasets
\newcommand{\ResSweepCostMultMax}{125}         % ... max over 7 datasets
\newcommand{\ResLadderBTwoCostMult}{2.0}      % ~ ladder-B=2 FLOPs / ours-B=1 FLOPs (typical)
\newcommand{\ResSweepVsBTwoMin}{5.4}            % GCCN-sweep FLOPs / ours-B=2 FLOPs, min
\newcommand{\ResSweepVsBTwoMax}{60.9}           % ... max
\newcommand{\ResCoraGap}{0.021}               % Cora: sweep-best minus ladder-B=2 test accuracy
\newcommand{\ResCiteseerGap}{0.013}           % Citeseer: sweep-best minus ladder-B=2 test accuracy
\newcommand{\ResFluorideBOneSd}{0.057}        % Fluoride-Carbonyl: ladder-B=1 test-accuracy sd (bimodal cascade)
\newcommand{\ResFluorideBTwoAcc}{0.939}       % Fluoride-Carbonyl: ladder-B=2 test accuracy (repair)
\newcommand{\ResFluorideBTwoSd}{0.006}        % Fluoride-Carbonyl: ladder-B=2 test-accuracy sd
\newcommand{\ResSelMaxSweepGap}{0.021}        % max over dataset rows of best-sweep minus best-selection accuracy (the table's tradeoff column)
\newcommand{\ResSelWithinRows}{6}            % dataset rows where the best selection arm matches/exceeds the best sweep within one sweep seed-sd
\newcommand{\ResSelTotalRows}{11}            % dataset rows with both a selection arm and a sweep

% ======================================================================
% Campaign-11 macros (11-dataset coverage, 2026-08-17): MUTAG live
% fixed-k tail + MANTRA cross-family contrast in balanced accuracy.
% Computed from the merged frozen batches by scripts/make_tables.py
% --performance-table (values below are patched by that script). UNFROZEN.
% ======================================================================
\newcommand{\ResLeadPanels}{11}               % datasets/panels in the lead figure (merged ladder-B=2 batches)
\newcommand{\ResMutagFixedKLiveAcc}{0.816}    % MUTAG: fixed anchored k=3, live test accuracy mean (anchored3_exact_t1)
\newcommand{\ResMutagFixedKLiveSd}{0.121}     % ... seed sd
\newcommand{\ResMutagLadderBTwoAcc}{0.915}    % MUTAG: ladder B=2 test accuracy mean (ladder_b2_t1)
\newcommand{\ResMutagLadderBTwoSd}{0.021}     % ... seed sd
\newcommand{\ResMutagExactOpt}{0.922}         % MUTAG: exact optimum of the enumerated 6-neighborhood landscape (#9 pooled)
\newcommand{\ResMantraGCCNLadderBal}{0.615}   % MANTRA: GCCN ladder B=2, BALANCED test accuracy mean
\newcommand{\ResMantraGCCNLadderBalSd}{0.103} % ... seed sd
\newcommand{\ResMantraAnchoredBal}{0.685}     % MANTRA: GCCN fixed anchored k=3, BALANCED test accuracy mean
\newcommand{\ResMantraAnchoredBalSd}{0.108}   % ... seed sd
\newcommand{\ResMantraHopseLadderBal}{0.771}  % MANTRA: HOPSE ladder B=2, BALANCED test accuracy mean
\newcommand{\ResMantraHopseLadderBalSd}{0.017}% ... seed sd
\newcommand{\ResMantraHopseKRungBal}{0.778}   % MANTRA: HOPSE fixed k=3 rung, BALANCED test accuracy mean
\newcommand{\ResMantraHopseKRungBalSd}{0.021} % ... seed sd
\newcommand{\ResMantraHopseSweepBal}{0.792}   % MANTRA: HOPSE sweep per-seed val-best, BALANCED test accuracy mean
\newcommand{\ResMantraHopseSweepBalSd}{0.009} % ... seed sd
\newcommand{\ResMantraSweepVsLadderMult}{7.2} % MANTRA: HOPSE-sweep grid FLOPs / HOPSE-ladder-B=2 FLOPs
\newcommand{\ResMantraMajorityRate}{0.921}    % MANTRA test-split majority-class rate (collapsed-classifier line)
\newcommand{\ResMantraCollapsedModels}{7}     % frozen MANTRA cells whose plain test accuracy is byte-identical to the majority rate

% --------------------------------------------------- protocol constants
\newcommand{\ResSeedList}{$\{42, 43, 44\}$}   % seeds shared by ours and every baseline arm
\newcommand{\ResOurReplications}{5}           % seed replications of our arms (supplementary table)
\newcommand{\ResCalibBackboneParams}{8704}    % backbone parameter count at the calibration mask (protocol-fidelity gate)
\newcommand{\ResCalibMaskID}{132}             % calibration mask (Benzene) used for the parameter match
\newcommand{\ResManifestHash}{715daee73cedf393}


% ----- working title (keep as a macro so it is easy to change) --------
\newcommand{\papertitle}{TopoSHAP: Game-Theoretic Explanations for
  Interpretable and Efficient Topological Deep Learning}

\title{\papertitle}
% Double-blind: the placeholder style prints the anonymous block below
% whenever \iclrfinalcopy is not set; \author is only used in final copy.
\author{Anonymous authors\\ Paper under double-blind review}

\begin{document}

\maketitle

\begin{abstract}
Topological deep learning (TDL) models interactions beyond pairwise
edges, such as rings in molecules and cavities in meshes. Existing
post-hoc explainers operate only on nodes and edges, so they cannot
represent the higher-order structures these models use. We introduce
\method, a model-agnostic method that explains predictions on
combinatorial complexes through a cooperative game over cells. It
produces per-cell Shapley attributions and identifies the smallest
sub-complex sufficient to preserve a prediction. For graphs, \method
reduces to induced-subgraph removal and matches or outperforms GNN
explainers on GraphXAI benchmarks while costing up to 500 times less
than search-based methods. The same implementation extends directly to
higher-order models: it recovers motifs in lifted molecular benchmarks,
achieves the highest faithfulness on topological datasets without
ground-truth explanations, and identifies decision-relevant structures
in native triangulations. Recasting neighborhoods as players and
validation accuracy as the game value also allows \method to select
architectures at roughly one-tenth the cost of a full sweep. Finally,
we find that standard feature liftings encode structural information in
features, reducing explainability. A structure-only lifting improves
both predictive accuracy and explanation quality.
\end{abstract}

% ======================================================================
% Section 1 -- Introduction (~1.25 pp incl. Fig 1)
% Binding: "what explainability buys practitioners"; "neighborhoods",
% never "relations"; no em dashes; short sentences; molecule gloss
% appears once, TopoTune parenthetical style.
% ======================================================================
\section{Introduction}
\label{sec:introduction}

Graph neural networks (GNNs) learn from pairwise connections. Many
systems carry structure beyond pairs: a benzene ring is a six-atom
object whose meaning is lost when read bond by bond. Topological deep
learning (TDL) models this higher-order structure directly
\citep{hajij2023tdl,papamarkou2024position}. Its domains are
\emph{combinatorial complexes}, which store multi-way interactions as
first-class objects called \emph{cells}, organized by \emph{rank}: in a
molecule, atoms (nodes, i.e., cells of rank zero), bonds (edges, i.e.,
cells of rank one), and rings (cells of rank two). Topological neural
networks (TNNs) exchange messages between cells along
\emph{neighborhoods} \citep{bodnar2021cwn,papillon2025topotune,
hopse2025}, and now match or exceed GNNs on molecular and geometric
benchmarks \citep{telyatnikov2024topobench}.

These models cannot yet be explained in their own vocabulary. Post-hoc
explainers for GNNs \citep{ying2019gnnexplainer,luo2020parameterized,
yuan2021subgraphx,muschalik2024graphshapiq} operate on nodes and
edges. A ring is not expressible in their output space, so the
structures a TNN reasons with cannot appear in its explanation. A
practitioner deploying a TNN today cannot answer two basic questions:
which cells drove this prediction, and which neighborhoods drive this
model's performance.

We introduce \method, which answers both questions with one cooperative
game (\cref{fig:overview}). The players are the cells of the input
complex. The value of a coalition is the model's output when only those
cells participate: features, neighborhood matrices, and readout pooling
are all restricted to the coalition. From this game \method returns two
outputs. The \emph{attribution} assigns each cell a Shapley value
\citep{shapley1953value}, inheriting the classical axioms. The
\emph{explanation} is the smallest sub-complex sufficient to preserve
the prediction, found by greedy pruning.

The game makes no assumption about rank. On graphs, which are rank-1
complexes, it reduces to induced-subgraph removal, and the explanation
output solves the same search problem as SubgraphX
\citep{yuan2021subgraphx}. This reduction lets us validate \method on
the community's own benchmarks before using its generality: on GraphXAI
\citep{agarwal2023graphxai}, \method matches or outperforms GNN
explainers on identical subjects, splits, and metrics
(\cref{tab:graph-benchmarks}), at up to 500 times lower cost than
search-based methods. The same implementation then runs unchanged on
higher-order models, where the baselines are undefined. It recovers
motifs on lifted molecular benchmarks (\cref{tab:lifted}), is the most
faithful method on TopoBench datasets that carry no explanation ground
truth (\cref{tab:nogt}), and separates how models decide on native
triangulations (\cref{sec:predictions}).

The same game template also explains performance
(\cref{sec:performance}). With neighborhoods as players and validation
accuracy as the value, one short training run prices every
message-passing route. The resulting selection recipe costs roughly one
tenth of an architecture sweep and concedes at most two accuracy
points to it.

Explaining these models also taught us something about their inputs.
Standard feature liftings compute higher-rank features by aggregating
lower-rank ones, which caches structural facts into feature vectors.
Models trained on such caches read them instead of computing structure,
and their explanations degrade measurably. A structure-only lifting
removes the cache and improves accuracy and explanation quality
together (\cref{sec:predictions}). This is a design rule practitioners
can apply directly.

\textbf{Contributions.}
\begin{itemize}
  \item \textbf{Method.} The participation game on combinatorial
    complexes, with two outputs: per-cell Shapley attributions and the
    smallest sufficient sub-complex (\cref{sec:framework}).
  \item \textbf{Reduction.} At rank 1 the game recovers the search
    problem of SubgraphX; the remaining design choices are ablated
    empirically (\cref{sec:predictions}).
  \item \textbf{Evidence.} Fair, same-subject comparisons on GraphXAI;
    motif recovery, ground-truth-free faithfulness, and model
    diagnostics on higher-order domains; all with multi-seed error
    bars (\cref{sec:predictions}).
  \item \textbf{Architecture selection.} The performance
    instantiation and its selection recipe, priced honestly
    (\cref{sec:performance}).
  \item \textbf{The lifting finding.} Feature caches reduce
    explainability; a structure-only lifting improves accuracy and
    explanations together (\cref{sec:predictions}).
  \item \textbf{Open source.} One reproducible pipeline covering every
    number and figure in the paper.
\end{itemize}

\begin{figure}[t]
  \centering
  \includegraphics[width=\linewidth]{../figures/out/overview.pdf}
  \caption{\textbf{\method: one game, two instantiations.}
  \textbf{A.}~Players are the cells of the input complex
  $\mathcal{C}$. A coalition $S$ keeps a subset of cells;
  $v(S)$ is the model output on the induced sub-complex
  $\mathcal{C}[S]$. Averaging marginal contributions over coalitions
  yields one Shapley value $\phi_\sigma$ per cell (output 1). Greedy
  pruning yields the smallest sufficient sub-complex $S^{*}_{b}$
  (output 2). \textbf{B.}~Players are the neighborhoods
  $\mathcal{N}$ of the architecture, drawn as source rank
  $\rightarrow$ destination rank, and $v(S)$ is the validation
  accuracy of a briefly trained model using only the neighborhoods in
  $S$; the same averaging attributes performance per neighborhood
  ($\phi_{\mathcal{N}}$).}
  \label{fig:overview}
\end{figure}

% ======================================================================
% Section 2 -- Related Work (~0.5 page budget)
% "neighborhoods", never "relations". No em dashes. Never a critique
% of TDL.
% ======================================================================
\section{Related Work}
\label{sec:related}

\paragraph{Explaining GNNs.}
Mask-optimization explainers learn continuous masks over edges and
features by gradient descent: GNNExplainer per instance
\citep{ying2019gnnexplainer} and PGExplainer amortized over a dataset
\citep{luo2020parameterized}. Search-based explainers return discrete
subgraphs scored by prediction preservation: SubgraphX explores
coalitions by Monte-Carlo tree search with sampled Shapley rewards
\citep{yuan2021subgraphx}, and Zorro greedily selects sparse masks with
stability guarantees \citep{funke2023zorro}. GraphSHAP-IQ computes
exact any-order Shapley interactions for GNNs
\citep{muschalik2024graphshapiq}. All of these operate on nodes and
edges. FORGE \citep{forge2025} lifts graphs to cell complexes before
explaining, yet still outputs node-level attributions. To our
knowledge, no prior method explains topological models natively at the
cell level. \Cref{sec:framework} makes the relationship precise: at
rank 1 our explanation output solves the search problem of SubgraphX,
and the remaining design differences are ablated in
\cref{sec:predictions}.

\paragraph{Topological deep learning.}
TDL generalizes message passing to simplicial complexes, cellular
complexes, hypergraphs, and combinatorial complexes
\citep{hajij2023tdl,papamarkou2024position}. Representative
architectures include CWN \citep{bodnar2021cwn}, SCCNN
\citep{yang2022sccnn}, CAN \citep{giusti2022can}, and CCXN
\citep{hajij2020ccxn}. TopoTune \citep{papillon2025topotune} defines
the GCCN design space we study: an architecture is specified by its
neighborhood set, which makes ``which neighborhoods matter'' a
first-class question. TopoBench \citep{telyatnikov2024topobench}
standardizes liftings, models, and training; all our experiments run on
it. HOPSE \citep{hopse2025} computes higher-order representations
without message passing. Our game treats HOPSE and message-passing
TNNs identically because it only requires forward evaluations.

\paragraph{Shapley values.}
The Shapley value is the unique attribution satisfying efficiency,
symmetry, null player, and linearity \citep{shapley1953value};
interaction indices extend it to coalitions
\citep{grabisch1999axiomatic}. SHAP popularized Shapley explanations
for machine-learning models \citep{lundberg2017shap}; permutation
sampling gives unbiased estimates at controlled cost
\citep{castro2009polynomial,strumbelj2010efficient}. Using Shapley
values to select players is unsound under redundancy
\citep{fryer2021shapley}. Our selection recipe therefore uses the game
only to order candidates and certifies every choice by validation over
trained models (\cref{sec:framework-game2}).

% ======================================================================
% Section 3 -- Preliminaries (~1 page budget)
% Say "neighborhoods", never "relations", for message-passing routes.
% ======================================================================
\section{Preliminaries}
\label{sec:preliminaries}

Three objects underlie the paper: the data structure (complexes and
cells), the models (GCCNs and neighborhoods), and the attribution
machinery (games and Shapley values). The molecule is the running
example: atoms are rank-0 cells, bonds rank-1, rings rank-2. A
glossary of every recurring term is in \cref{app:glossary}.

\subsection{Combinatorial Complexes}
\label{sec:prelim-cc}

\begin{definition}[Combinatorial complex \citep{hajij2023tdl}]
\label{def:cc}
A \emph{combinatorial complex} is a triple $(V, \mathcal{X}, \rk)$ where
$V$ is a finite set of vertices, $\mathcal{X} \subseteq
2^{V}\setminus\{\emptyset\}$ is a set of \emph{cells}, and $\rk :
\mathcal{X} \to \mathbb{Z}_{\ge 0}$ is a rank function monotone with
respect to inclusion: $x \subseteq y \implies \rk(x) \le \rk(y)$.
\end{definition}

Cells of rank $0$ and $1$ are nodes and edges; cells of rank $2$ are
faces, and higher ranks hold further multi-way structure. Graphs are
the rank-$\le 1$ special case. Hypergraphs carry set-type relations
and no rank-2 cells, so nothing in what follows assumes rank-2 cells
exist. A \emph{lifting} converts a graph into a complex by promoting
recognizable substructures to higher-rank cells, \eg chordless cycles
to faces \citep{bodnar2021cwn}; MANTRA
\citep{ballester2025mantra} is natively higher-order and needs no
lifting. We write $\mathcal{X}_r$ for the cells of rank $r$ and
$\mathbf{H}_r$ for the feature matrix on rank $r$.

\subsection{GCCNs and Neighborhoods}
\label{sec:prelim-gccn}

Message passing on a complex is organized by \emph{neighborhoods}:
routes along which cells exchange messages, each named by who talks to
whom \citep{papillon2025topotune}: \emph{incidences} $B_{r,r'}$
(\eg nodes to the edges containing them), \emph{up-adjacencies}
$A^{\uparrow}_{r}$ (\eg nodes joined by an edge, written
\nbhd{up\_adjacency-0}), and \emph{down-adjacencies}
$A^{\downarrow}_{r}$. A generalized combinatorial complex network
(GCCN) is specified by a set $\mathcal{N} = \{N_1, \dots, N_m\}$ of
neighborhoods; each layer runs one message-passing submodule per
neighborhood and aggregates:
\begin{equation}
\label{eq:gccn-layer}
\mathbf{h}^{(\ell+1)}_x
  \;=\;
  \phi\!\Bigl(\mathbf{h}^{(\ell)}_x,\;
  \bigotimes_{N \in \mathcal{N}}\;
  \bigoplus_{y \in N(x)}
  \psi_{N}\bigl(\mathbf{h}^{(\ell)}_x, \mathbf{h}^{(\ell)}_y\bigr)\Bigr),
\end{equation}
with $\bigoplus$ aggregating within and $\bigotimes$ across
neighborhoods. In TopoBench the candidate set for rank-$\le 2$
complexes has \unfrozen{\ResNumNeighborhoods} neighborhoods; which
subset to activate is ordinarily answered by sweeps that train one
model per candidate subset \citep{papillon2025topotune}.
\Cref{sec:performance} answers it with one training run of the
all-neighborhoods model plus attribution.

\subsection{Shapley Values and Interactions}
\label{sec:prelim-shapley}

A \emph{cooperative game} $(P, v)$ has players $P$, $|P| = n$, and a
value function $v : 2^{P} \to \mathbb{R}$, $v(\emptyset) = 0$, scoring
every \emph{coalition} (subset). The \emph{Shapley value}
\citep{shapley1953value} of player $i$ averages its marginal
contribution over all coalitions:
\begin{equation}
\label{eq:shapley}
\phi_i(v) \;=\; \sum_{S \subseteq P \setminus \{i\}}
  \frac{|S|!\,(n - |S| - 1)!}{n!}
  \bigl[v(S \cup \{i\}) - v(S)\bigr].
\end{equation}
It is the unique attribution satisfying \emph{efficiency}
($\sum_i \phi_i = v(P)$), \emph{symmetry}, \emph{null player}
($\phi_i = 0$ if $i$ never changes any coalition's value), and
\emph{linearity}. The \emph{interaction index}
\citep{grabisch1999axiomatic} extends \cref{eq:shapley} to coalitions
$T$ via discrete derivatives:
\begin{equation}
\label{eq:interaction}
I(T) \;=\; \sum_{S \subseteq P \setminus T}
  \frac{|S|!\,(n - |S| - |T|)!}{(n - |T| + 1)!}\,
  \Delta_{T} v(S),
\qquad
\Delta_{T} v(S) = \sum_{L \subseteq T} (-1)^{|T| - |L|}\, v(S \cup L),
\end{equation}
so $I(\{i,j\}) > 0$ indicates synergy and $I(\{i,j\}) < 0$ redundancy.
Exact computation needs $2^{n}$ evaluations; \cref{sec:framework}
gives the regimes where our games afford it and the sampled estimator
otherwise. Because the axioms are equalities a finite computation can
check, they double as executable tests of the model-plus-explainer
stack (\cref{app:experimental}).

% ======================================================================
% Section 4 -- METHOD (~2 page budget). Direct register, no em dashes.
% ======================================================================
\section{Method: One Game, Two Outputs}
\label{sec:framework}

\method plays one cooperative game and reads two answers from it. The
players are the cells of the input complex. The value of a coalition is
the model's output when only those cells participate. The game returns
\emph{attributions} (per-cell Shapley values, \cref{eq:shapley}) and an
\emph{explanation} (the smallest sufficient sub-complex). The same
template, with neighborhoods as players and validation accuracy as the
value, explains \emph{performance} (\cref{sec:framework-game2}).

\subsection{The Participation Game}
\label{sec:framework-game1}

To explain a prediction $f(\mathcal{C})$ for class $y$, the players are
the cells $P \subseteq \mathcal{X}$ of the input, by default all cells,
with $n = |P|$. A coalition $S$ is scored on the \emph{induced
sub-complex} $\mathcal{C}[S]$: cells outside $S$ are detached, meaning
their features are removed, every neighborhood matrix is recomputed
from the restricted incidences, and readout pooling is limited to $S$:
\begin{equation}
\label{eq:participation-game}
v(S) \;=\;
  f_{y}\bigl(\mathcal{C}[S]\bigr) - f_{y}\bigl(\mathcal{C}[\emptyset]\bigr).
\end{equation}
Detachment is what absence means for a cell. Zeroing features instead
leaves the cell in place, still passing messages and still contributing
bias terms to pooling, so a zeroed game scores a different
counterfactual (quantified in \cref{app:experimental}). Detachment is
also cell-native: removing a face removes the rank-2 player itself.

\paragraph{Output 1: attributions.}
Shapley values $\phi_\sigma(v)$ on \cref{eq:participation-game} give
signed per-cell credit with the axioms of \cref{sec:prelim-shapley}.
Efficiency decomposes the prediction over cells, and the null-player
axiom detects models that ignore their higher-order structure
(\cref{sec:predictions}). Small complexes are enumerated exactly.
Larger ones use $T$ permutation passes, which cost $T \cdot n$ game
evaluations and are unbiased \citep{castro2009polynomial}.

\paragraph{Output 2: explanations.}
Motif recovery asks for a sub-complex, so \method extracts one
directly: the budget-$b$ coalition
$S^{*}_{b} \approx \argmax_{|S| = b} v(S)$, found by greedy pruning.
We start from $P$, repeatedly detach the cell whose removal least
reduces $v$, and record $S$ at every requested budget. One pass costs
at most $(n^{2}+n)/2$ evaluations of one forward pass each and yields
the coalition at all budgets, since pruning is nested. Pruning from
the full complex keeps evaluations near the training distribution.
Growing coalitions from single cells scores far outside it and fails
empirically (\cref{app:experimental}).

\paragraph{Relation to graph explainers.}
A graph is a complex of rank $\le 1$, where
\cref{eq:participation-game} reduces to induced-subgraph node removal.
Output 2 then solves the same search problem as SubgraphX
\citep{yuan2021subgraphx}, with two substitutions: detachment replaces
feature zero-filling, and exhaustive one-step pruning, at one forward
pass per scored coalition, replaces Monte-Carlo tree search whose
reward samples ${\sim}10^{2}$ context coalitions per node it scores.
Both substitutions are ablated in \cref{sec:predictions}. GNNExplainer
\citep{ying2019gnnexplainer} and PGExplainer
\citep{luo2020parameterized} optimize continuous relaxed masks and
return no game or axioms. All three are defined only at rank $\le 1$.

\subsection{Neighborhoods Explain Performance}
\label{sec:framework-game2}
\label{sec:method-pipeline}

The same template explains performance. The players are the model's
neighborhood-indexed components: a GCCN has one message-passing route
per neighborhood and HOPSE has one encoding block per neighborhood, so
switching a player off removes that neighborhood. We propose two value
functions. The \emph{retraining value} $v_{\mathrm{retrain}}(S)$ is
the validation accuracy of the model trained from scratch with the
neighborhoods in $S$; it is exact, and each evaluation costs a full
training run. The \emph{masked value} $v_{\mathrm{mask}}(S)$ is the
validation accuracy of a briefly trained full model evaluated with the
neighborhoods outside $S$ removed; each evaluation costs one forward
pass over the validation set.

We turn this game into an architecture selection procedure:
\begin{recipebox}
  \item We train the full model, with all candidate neighborhoods, for
    a small fixed number of epochs
    ($\sim$\unfrozen{\ResStemEpochs}).
  \item We repeatedly remove the neighborhood whose removal costs the
    least masked value, while that cost stays below the validation
    noise floor. This yields one candidate architecture per size.
  \item We train $B$ of these candidates to completion and keep the
    one with the best validation accuracy.
\end{recipebox}

The stop threshold is the binomial noise floor $\sqrt{p(1-p)/m}$ of a
validation set with $m$ examples, known before any training, so the
selected size is not a hyperparameter. The budget $B$ trades cost
against reliability, from $\sim$\unfrozen{\ResLadderCostBOne} training
runs at $B{=}1$ to $\sim$\unfrozen{\ResLadderCostBTwo} at $B{=}2$; a
full sweep trains every size. Two rules protect the procedure. The
masked value only orders candidates, and every accept or reject
decision is validation of fully trained models. Elimination walks
downward from the full model, the regime where the masked value is
most faithful. Details and supporting evidence are in
\cref{app:selection}.

\subsection{Exactness}
\label{sec:framework-exact}
\label{sec:framework-theory}

The neighborhood game has $n = \unfrozen{\ResNumNeighborhoods}$
players and $2^{n} = 2048$ coalitions, so \cref{eq:shapley,eq:interaction}
are evaluated exactly. Prediction games on small complexes likewise
enumerate; larger instances use sampled passes. Against brute force,
the implementation agrees to \unfrozen{\ResExactnessMaxErr}, which is
floating-point exactness. Exact axioms double as executable software
checks; the bugs they caught are documented in
\cref{app:experimental}.

\begin{theorybox}
\begin{propositionT}[Non-identifiability of substitutes; statement placeholder]
% (no \label: unnumbered T-statement; refer to it as "Proposition~T1")
% TODO(researcher): finalize statement + proof.
Let $N_a, N_b \in \mathcal{N}$ be neighborhoods that are functionally
substitutable for a trained model $f_\theta$, i.e., there exist
parameters realizing the same function with either route active. Then
the masked value computed on any single trained model cannot
distinguish the attribution profile of $(N_a, N_b)$ from that of a
model committed to either substitute. \emph{Exact statement TBD by
researcher.}
\end{propositionT}
\end{theorybox}

Proposition~T1 states which game answers which question: the masked
value identifies the trained model's realized division of labor, and
claims about what is realizable under retraining require the
retraining value. The selection procedure operationalizes this by
letting validation over trained models decide.

% ======================================================================
% Section 5 -- Results I: Explaining predictions (~2.4 pp)
% Bold-first-sentence paragraphs (\expclaim). Numbers live in tables.
% No em dashes. Multi-seed error bars everywhere.
% ======================================================================
\section{Results I: Explaining Predictions}
\label{sec:predictions}

\textbf{Setup.} We use the four GraphXAI molecular benchmarks with
annotated ground-truth motifs \citep{agarwal2023graphxai}, the
TopoBench datasets NCI1 and PROTEINS, which carry no explanation
annotations, and MANTRA \citep{ballester2025mantra}, a native
simplicial dataset. Every comparison uses identical training splits,
the full ground-truth-positive test population, and three training
seeds; every entry reports mean $\pm$ standard deviation over seeds.
These experiments realize the framework on
cellular complexes (the molecular benchmarks, lifted) and simplicial
complexes (MANTRA). Explanation accuracy is measured by GEA: select $b$ cells, where $b$ is
the size of a ground-truth candidate, and report the Jaccard overlap
with that candidate, maximized over candidates. Faithfulness is
measured by deletion-AUC: detach cells in score order and integrate the
predicted-class probability, so lower is better. Full protocols are in
\cref{app:experimental}.

\subsection{Validation at Rank 1: Graph Benchmarks}
\label{sec:predictions-rank1}

On graphs the participation game reduces to induced-subgraph removal
(\cref{sec:framework}). This lets us compare \method against GNN
explainers on their own terms. All methods explain the same trained,
parameter-matched GCN and GIN subjects.

% AUTO-GENERATED by scripts/make_paper_tables.py -- do not edit
\begin{table}[t]
\centering
\small
\caption{\textbf{Graph benchmarks (rank-1 validation).} Node-level explanation accuracy (GEA, mean $\pm$ sd over 3 training seeds) on the full ground-truth-positive test splits. All methods explain the same parameter-matched GNN subjects on identical splits; per method, the better of the GCN/GIN subjects is shown (both in \cref{app:tables}). \method{} explanations cost $\sim$0.04\,s per instance vs.\ $\sim$20\,s for SubgraphX on identical hardware.}
\label{tab:graph-benchmarks}
\begin{tabular}{lcccc}
\toprule
 & Benzene & FluorideCarbonyl & Mutagenicity & AlkaneCarbonyl \\
\midrule
\method{} (attribution) & 0.859 $\pm$ 0.009 & \cellcolor{black!10}0.507 $\pm$ 0.047 & \cellcolor{black!10}\textbf{0.798 $\pm$ 0.124} & \cellcolor{black!10}\textbf{0.500 $\pm$ 0.081} \\
\method{} (explanation) & \cellcolor{black!10}\textbf{0.969 $\pm$ 0.009} & 0.340 $\pm$ 0.011 & 0.480 $\pm$ 0.135 & 0.354 $\pm$ 0.125 \\
SubgraphX & 0.796 $\pm$ 0.029 & 0.439 $\pm$ 0.107 & 0.344 $\pm$ 0.036 & 0.296 $\pm$ 0.010 \\
PGExplainer & 0.642 $\pm$ 0.044 & \cellcolor{black!10}\textbf{0.555 $\pm$ 0.067} & 0.290 $\pm$ 0.102 & 0.205 $\pm$ 0.160 \\
GNNExplainer & 0.443 $\pm$ 0.054 & 0.210 $\pm$ 0.014 & 0.158 $\pm$ 0.023 & 0.162 $\pm$ 0.049 \\
Random & 0.250 $\pm$ 0.000 & 0.219 $\pm$ 0.000 & 0.132 $\pm$ 0.000 & 0.173 $\pm$ 0.000 \\
\bottomrule
\end{tabular}
\end{table}


\expclaim{\method matches or outperforms every baseline on the
GraphXAI benchmarks.} \Cref{tab:graph-benchmarks} reports both
outputs. The explanation output leads on Benzene and the attribution
output leads on Mutagenicity and Alkane-Carbonyl. PGExplainer retains
the best mean on Fluoride-Carbonyl, within overlapping error bars of
our attribution output. The two outputs are complementary: attribution
measures how much the prediction depends on each cell, and wins where
the motif is necessary; explanation measures joint sufficiency, and
wins where the motif alone preserves the prediction.

\expclaim{The explanation output is orders of magnitude cheaper than
search-based explanation.} Greedy pruning evaluates each candidate
removal once per step, one forward pass per coalition. SubgraphX pays
roughly two hundred forward passes per coalition it scores, because its
Monte-Carlo reward samples coalition contexts. Measured wall-clock
costs on identical hardware appear in the caption of
\cref{tab:graph-benchmarks}.

\expclaim{Both removal semantics and search direction matter, and both
choices are ablated.} Replacing detachment with feature zero-filling,
the convention of SubgraphX, reproduces SubgraphX's behavior on GCN
subjects and degrades substantially on GIN subjects
(\cref{tab:by-subject}). Growing coalitions from single cells instead of
pruning from the full complex fails: early evaluations score
coalitions far outside the training distribution
(\cref{app:experimental}).

\expclaim{Where a baseline wins, the game explains why.} Define the
alignment gap $\Delta = v(S^{*}_{b}) - v(S_{\mathrm{gt}})$, where
$S_{\mathrm{gt}}$ is the ground-truth motif as a coalition. On
Fluoride-Carbonyl the explained GIN scores non-motif coalitions above
the annotated motif on 96\% of candidates, so every faithful explainer
is capped there. GEA evaluates the composition of explainer and model
against human annotations. $\Delta$ separates the two factors at the
cost of two extra game evaluations per instance, and no baseline can
compute it.

\subsection{Higher-Order Domains}
\label{sec:predictions-higher}

We now explain TNNs on lifted complexes, where faces are rank-2
players. The explained model in every experiment is the
full-vocabulary GCCN, which uses all eleven candidate neighborhoods
and involves no selection step.

\expclaim{Standard feature liftings cache structure into features, and
models trained on them are harder to explain.} The standard lifting
computes a face's features by summing its nodes' features at
preprocessing time. The face then carries information about its
boundary even after the boundary is detached, so the counterfactuals
of any removal-based explainer leak. Trained models exploit the cache:
their sufficient coalitions keep the face while dropping its
boundary (\cref{app:experimental}), and their alignment gap is
positive on nearly all candidates.

% AUTO-GENERATED by scripts/make_paper_tables.py -- do not edit
\begin{table}[t]
\centering
\small
\caption{\textbf{Lifted topological benchmarks.} Node-projected GEA (mean $\pm$ sd, 3 seeds, full populations) for both \method{} outputs on the same full-vocabulary GCCN trained under two feature liftings. The standard lifting sums boundary features into higher-rank cells; the structural lifting uses constant features. Bold marks the best entry per column. Subject accuracies: Benzene 0.942, FluorideCarbonyl 0.947, Mutagenicity 0.976, AlkaneCarbonyl 0.967; alignment-ceiling share $\Pr[v(S^{*})>v(S_{\mathrm{gt}})]$: Benzene 0.74, FluorideCarbonyl 0.93, Mutagenicity 0.97, AlkaneCarbonyl 0.99.}
\label{tab:lifted}
\begin{tabular}{lcccc}
\toprule
 & Benzene & FluorideCarbonyl & Mutagenicity & AlkaneCarbonyl \\
\midrule
attribution, structural lifting & 0.471 $\pm$ 0.018 & \cellcolor{black!10}0.491 $\pm$ 0.035 & 0.326 $\pm$ 0.065 & \cellcolor{black!10}0.333 $\pm$ 0.097 \\
attribution, standard lifting & 0.459 $\pm$ 0.192 & \cellcolor{black!10}\textbf{0.516 $\pm$ 0.089} & 0.257 $\pm$ 0.077 & \cellcolor{black!10}0.348 $\pm$ 0.043 \\
explanation, structural lifting & \cellcolor{black!10}\textbf{0.725 $\pm$ 0.105} & 0.426 $\pm$ 0.027 & \cellcolor{black!10}\textbf{0.436 $\pm$ 0.081} & \cellcolor{black!10}0.362 $\pm$ 0.130 \\
explanation, standard lifting & 0.509 $\pm$ 0.188 & 0.372 $\pm$ 0.004 & 0.289 $\pm$ 0.117 & \cellcolor{black!10}\textbf{0.364 $\pm$ 0.053} \\
\bottomrule
\end{tabular}
\end{table}


\expclaim{A structure-only lifting improves accuracy and explanation
quality together.} The structural lifting assigns constant features to
all cells above rank 0, the convention MANTRA uses natively, so
structure exists only in the incidences and must be computed by
message passing. \Cref{tab:lifted} compares the same architecture
under both liftings: the structural lifting raises GEA, lowers the
alignment-ceiling share, and trains to higher accuracy. Explanations
of the structural models recover the face's boundary, and the face
itself appears in the explanation on 92\% of Benzene instances.
A cell-level pointer of this kind is not expressible by any node-based
method.

% AUTO-GENERATED by scripts/make_paper_tables.py -- do not edit
\begin{table}[t]
\centering
\small
\caption{\textbf{Faithfulness without ground truth.} Deletion-AUC (lower is better; predicted-class probability as cells are detached in score order; mean $\pm$ sd over 3 training seeds, 50 pre-declared test complexes each) on TopoBench datasets with no explanation annotations. Subjects: full-vocabulary GCCNs over the structural lift; explanations target the model's own prediction.}
\label{tab:nogt}
\begin{tabular}{lcccc}
\toprule
 & NCI1 & NCI109 & MUTAG & PROTEINS \\
\midrule
\method{} & \cellcolor{black!10}\textbf{0.218 $\pm$ 0.005} & \cellcolor{black!10}\textbf{0.203 $\pm$ 0.044} & \cellcolor{black!10}\textbf{0.175 $\pm$ 0.043} & \cellcolor{black!10}\textbf{0.338 $\pm$ 0.067} \\
Occlusion & 0.254 $\pm$ 0.021 & \cellcolor{black!10}0.240 $\pm$ 0.038 & 0.231 $\pm$ 0.006 & 0.431 $\pm$ 0.065 \\
Gradient$\times$Input & 0.552 $\pm$ 0.016 & 0.507 $\pm$ 0.049 & 0.499 $\pm$ 0.073 & 0.638 $\pm$ 0.047 \\
Random & 0.545 $\pm$ 0.004 & 0.491 $\pm$ 0.035 & 0.481 $\pm$ 0.014 & 0.666 $\pm$ 0.058 \\
\bottomrule
\end{tabular}
\end{table}


\expclaim{Without ground truth, \method is the most faithful method
and reports how much topology a model uses.} On NCI1 and PROTEINS
there are no annotated motifs, which is the common case in practice.
Explanations target the model's own prediction
(\cref{sec:framework-game1}). \Cref{tab:nogt} shows deletion-AUC
against occlusion, gradient-times-input, and random baselines on the
same pre-declared test samples. The attribution output also yields a
diagnostic no baseline provides: the share of Shapley mass per rank.
On NCI1 the trained models place roughly half of their credit on
rank-1 cells and less than a tenth on rank-2 cells, a direct readout
of how much higher-order structure each model actually uses.

\expclaim{On native triangulations, rank-level attributions separate
how models decide.} On MANTRA's orientability task, per-rank
attribution profiles distinguish model families that reach similar
accuracy through different mechanisms, and the null-player axiom
exposed one published configuration whose higher-rank cells received
exactly zero credit: a topological model that never used its topology.
\method is also the most faithful method on MANTRA under deletion-AUC.
Full profiles and the multi-complex protocol are in
\cref{app:figures,app:experimental}.

% ======================================================================
% Section 6 -- Results II: Explaining performance (~1.5 pp).
% Direct register; no em dashes; no "ladder/rung/stem" vocabulary.
% BINDING: multi-seed error bars; parameter/budget matching; linear
% axes; uniform y-span; one star per fixed k; menu at true cost.
% ======================================================================
\section{Results II: Explaining Performance}
\label{sec:performance}

The selection procedure of \cref{sec:framework-game2} runs identically
on all \unfrozen{\ResLeadPanels} datasets and in two model families,
GCCN message-passing routes and HOPSE encoding blocks. We compare
against each family's own architecture sweep, with every method priced
by its total training compute in FLOPs under one measurement
convention (\cref{app:flops-conventions}). All markers carry
multi-seed error bars over identical seeds. \Cref{tab:performance-summary}
(\cref{app:selection-table}) tabulates every marker with exact values,
the $B{=}1$ arm, and per-row cost multiples. Comparisons against
published fixed architectures and fixed-$k$ variants are in
\cref{app:selection}.

\begin{figure}[!t]
  \centering
  \IfFileExists{../figures/out/lead_xai_anchored.pdf}{%
    \includegraphics[width=\linewidth]{../figures/out/lead_xai_anchored.pdf}%
  }{%
    \figureplaceholder{figures/out/lead_xai_anchored.pdf}%
  }
  \caption{\textbf{Architecture selection on the GraphXAI datasets.}
  Test accuracy vs.\ total training compute (FLOPs, linear axes).
  Dark markers are the architecture chosen by our selection procedure
  in each family; pale markers are each family's sweep at its full
  grid cost. Bars are sd over seeds.}
  \label{fig:lead}
\end{figure}

\begin{figure}[!t]
  \centering
  \IfFileExists{../figures/out/lead_bench_anchored.pdf}{%
    \includegraphics[width=\linewidth]{../figures/out/lead_bench_anchored.pdf}%
  }{%
    \figureplaceholder{figures/out/lead_bench_anchored.pdf}%
  }
  \caption{\textbf{Architecture selection on the TopoBench
  benchmarks.} Same conventions as \cref{fig:lead}. Top row:
  complex-level tasks; bottom row: node-level citation complexes. The
  HOPSE family is evaluated on complex-level tasks, so the node-level
  panels show the GCCN family. On MUTAG and NCI109 the GCCN sweep
  marker is the best of the 64-configuration retraining enumeration
  at estimated cost. MANTRA is scored in balanced accuracy (majority
  class \unfrozen{\ResMantraMajorityRate}).}
  \label{fig:lead-bench}
\end{figure}

\expclaim{Selection concedes at most \unfrozen{\ResSelMaxSweepGap}
accuracy for a
\unfrozen{\ResSweepCostMultMin}--\unfrozen{\ResSweepCostMultMax}$\times$
compute reduction.} In \cref{tab:performance-summary}, the best
selection arm matches or exceeds the best sweep within one sweep
seed-sd on \unfrozen{\ResSelWithinRows} of
\unfrozen{\ResSelTotalRows} datasets. On the remaining datasets the
sweep ends higher, never by more than \unfrozen{\ResSelMaxSweepGap}.
The largest concessions sit where the masked game carries the least
signal: the node-level citation complexes, whose small validation
splits raise the noise floor the stop rule reads (next claim), the
class-imbalanced MANTRA, and Alkane-Carbonyl, where the comparator is
the exhaustive enumeration rather than a sweep. The
$B{=}2$ budget costs
$\sim$\unfrozen{\ResLadderBTwoCostMult}$\times$ more, keeps a
\unfrozen{\ResSweepVsBTwoMin}--\unfrozen{\ResSweepVsBTwoMax}$\times$
compute margin over the sweep, and stabilizes the one dataset where
$B{=}1$ is noisy (Fluoride-Carbonyl,
\cref{tab:performance-summary}). The same procedure, unchanged, runs
in the HOPSE family (\cref{fig:lead,fig:lead-bench}). These
costs use the masked value; evaluating the retraining value instead
costs \unfrozen{\ResSweepEqMin}--\unfrozen{\ResSweepEqMax}
sweep-equivalents, so the masked value carries every cost claim
(regret $\le$ \unfrozen{\ResRegretMax} against retraining-value
selections, \cref{app:accounting}).

\expclaim{The number of neighborhoods kept is not a hyperparameter.}
The stop derives from the validation noise floor $\sqrt{p(1-p)/m}$,
which matches measured retraining noise in order of magnitude on all
\unfrozen{\ResBinomOrderDatasets} datasets. The rule was frozen before
its final scoring. Out of sample, the sequence of candidates contains
the exact optimum on
\unfrozen{\ResLadderOOSContained}/\unfrozen{\ResLadderOOSTotal}
never-touched architecture landscapes, with worst regret
\unfrozen{\ResLadderOOSWorstRegret}. Fixing $k$ in advance is the
alternative, and it fails where the right size differs: on MUTAG the
fixed $k{=}3$ model collapses while the selection holds near the
exact optimum (\cref{tab:performance-summary}).

\expclaim{Cross-family gaps are model-class facts that no
within-family selection can close.} On Benzene, HOPSE's spectral
encodings clear the ceiling of the entire GCCN neighborhood space by
$\sim$\unfrozen{\ResBenzeneSpectralMargin}. On MANTRA orientability,
the GCCN family sits near the count-only score of
\unfrozen{\ResCountingCeiling} in balanced accuracy while the same
procedure in the HOPSE family clears it at
$\sim$\unfrozen{\ResMantraSweepVsLadderMult}$\times$ less compute
than the HOPSE sweep (\cref{fig:lead-bench}). The GCCN sweep makes
the same point at full grid cost: its validation-best configuration
is a majority collapse (\cref{tab:performance-summary}). Plain-accuracy MANTRA
numbers collapse to the majority class, so all MANTRA results are
balanced accuracy (\cref{app:experimental}).


% ======================================================================
% Section 8 -- Limitations (one paragraph, concise). No em dashes.
% ======================================================================
\section{Limitations}
\label{sec:limitations}

\method has three main limitations. First, although greedy pruning is
far cheaper than search-based explainers at benchmark scale
(\cref{tab:graph-benchmarks}), its cost grows quadratically with the
number of cells, so complexes with thousands of cells would require
restricting the search to the model's receptive field. Second, GEA
scores the composition of explainer and model against human
annotations. On Fluoride-Carbonyl, for example, the explained model
scores coalitions outside the annotated functional group above the
group itself on 96\% of candidates, so no faithful method can reach a
high GEA there. We report the alignment gap $\Delta$ alongside GEA to
separate explainer quality from this model property. Third, all
prediction-explanation experiments target complex-level
classification; explaining node-level predictions would restrict
players to the target's receptive field, and we leave it to future
work.

% ======================================================================
% Section 9 -- Conclusion (~0.2 page budget). No em dashes.
% ======================================================================
\section{Conclusion}
\label{sec:conclusion}

\method explains topological deep learning with one cooperative game
on the cells of a combinatorial complex. The game returns per-cell
Shapley attributions and the smallest sufficient sub-complex. At rank
1 it matches or outperforms GNN explainers on their own benchmarks at
a small fraction of their cost. On higher-order domains, where those
explainers are undefined, it recovers motifs, leads on faithfulness
without ground truth, and reads out how much topology a model uses.
With neighborhoods as players, the same game selects architectures at
roughly one tenth of a sweep's cost. Its diagnostics produced a
concrete design rule: keep structural information out of lifted
features, and both accuracy and explainability improve. We release the
full pipeline, so every number and figure in this paper regenerates
from one command.
% TODO(researcher): anonymized repo link for submission.


% ----- required/recommended statements (do not count toward page limit)
% TODO(researcher): review wording of both statements before submission.
\subsection*{AI use statement}
We used generative AI tools (Claude) for code implementation, experiment
orchestration, data analysis, and manuscript drafting, in every case
under author direction and review. The research questions, method
design, experimental protocol, and framing are the authors' own. All
AI-assisted code is exercised by the released test suite, the exactness
of the attribution engine is verified against brute-force enumeration
to machine precision, and every number and figure in the paper
regenerates from the frozen experimental record by one command. We take
responsibility for the final content of this work.

\subsection*{Reproducibility statement}
Every number and figure in this paper is generated from a frozen
experimental record by a single script and referenced in the text as a
macro; none is typed by hand. The exact selection protocol, seeds,
FLOPs conventions, and dataset scope are specified in
\cref{app:experimental}, and the evidence layers (frozen rule,
out-of-sample validation, pre-registered battery) in
\cref{app:selection}. The full pipeline, including runners, frozen
results, tests, and the paper build, is released as open source upon
publication.

% ----- bibliography ----------------------------------------------------
% TODO: switch to the official iclr2027_conference.bst when released.
\bibliographystyle{plainnat}
\bibliography{references}

\appendix
% ======================================================================
% Appendix -- Notation and glossary
% Accessibility requirement (researcher-set): every recurring term and
% symbol, its plain-English meaning, and where it is first used.
% ======================================================================
\section{Notation and Glossary}
\label{app:glossary}

\Cref{tab:glossary} collects the paper's recurring terms and symbols.
Each entry gives the plain-English meaning used throughout, anchored in
the molecule running example (atoms $=$ nodes $=$ rank 0, bonds $=$
edges $=$ rank 1, rings $=$ faces $=$ rank 2), and the section where
the term is first defined.

\begin{table}[h]
\centering
\caption{Glossary of terms and symbols.}
\label{tab:glossary}
\small
\begin{tabular}{p{0.20\linewidth}p{0.56\linewidth}l}
\toprule
Term / symbol & Plain meaning & First use \\
\midrule
cell & A node or a group of nodes stored as one object: an atom, a
  bond, a ring. & \cref{sec:prelim-cc} \\
rank, $\rk$ & Which level a cell lives on: atoms 0, bonds 1, rings 2.
  & \cref{sec:prelim-cc} \\
combinatorial complex, $\mathcal{C}$ & The data structure holding all
  cells of all ranks for one input. & \cref{sec:prelim-cc} \\
lifting & Converting a plain graph into a complex, \eg promoting a
  molecule's chordless cycles to ring cells. & \cref{sec:prelim-cc} \\
neighborhood, $N$ & A message-passing route: who talks to whom (\eg
  atom$\to$bond, bond$\to$ring). & \cref{sec:prelim-gccn} \\
GCCN & Generalized combinatorial complex network: a TDL model built
  from one submodule per neighborhood. & \cref{sec:prelim-gccn} \\
neighborhood-indexed component & The architecture piece a model has
  one of per neighborhood: a GCCN route, a HOPSE encoding block.
  & \cref{sec:framework-game2} \\
game, $v$ & An assignment of a score to every subset of players.
  & \cref{sec:prelim-shapley} \\
coalition, $S$ & A subset of players (cells, or components).
  & \cref{sec:prelim-shapley} \\
marginal contribution & How much the score changes when one player
  joins a coalition. & \cref{sec:prelim-shapley} \\
Shapley value, $\phi_i$ & A player's average marginal contribution over
  all coalitions; the unique fair credit split.
  & \cref{sec:prelim-shapley} \\
masking & Neutralizing the players outside a coalition so the model
  behaves as if they were absent. & \cref{sec:framework-game1} \\
exact (retraining) game & Score of a coalition $=$ validation accuracy
  after training from scratch with exactly those components.
  & \cref{sec:framework-game2} \\
cheap (masked-validation) game, $v_{\mathrm{prune}}$ & Score of a
  coalition $=$ the brief run's validation accuracy with the other
  components removed; one forward pass per coalition.
  & \cref{sec:framework-game2} \\
brief run & The all-components model trained briefly
  (\unfrozen{\ResStemEpochs} epochs) before attribution.
  & \cref{sec:framework-game2} \\
candidate sequence & The nested sequence of shrinking architectures
  produced by backward elimination; one candidate per size.
  & \cref{sec:method-pipeline} \\
candidate & One architecture in the candidate sequence.
  & \cref{sec:method-pipeline} \\
stop candidate & The smallest candidate whose next removal would cost
  more than the noise floor. & \cref{app:selection-ladder} \\
budget $B$ & How many candidates are trained to completion before
  validation picks; the practitioner's dial.
  & \cref{sec:method-pipeline} \\
binomial floor & The validation set's own noise level,
  $\sqrt{p(1-p)/m}$; the elimination's stop threshold.
  & \cref{sec:method-pipeline} \\
substitute-commitment & The brief run's arbitrary choice between two
  components that can do the same job; attribution reports the choice,
  not the interchangeability. & \cref{app:selection-commitment} \\
menu & The \unfrozen{\ResMenuEntries} published fixed architectures
  from the literature, treated as one baseline at its true cost.
  & \cref{sec:method-pipeline} \\
sweep & Training one model per candidate configuration and keeping the
  best; the procedure the pipeline replaces. & \cref{sec:prelim-gccn} \\
\bottomrule
\end{tabular}
\end{table}

% ======================================================================
% Appendix A -- Proofs
% ======================================================================
\section{Proofs}
\label{app:proofs}

\subsection{Proof of Proposition T1 (Non-Identifiability of Substitutes)}
\label{app:proof-t1}

% TODO(researcher): finalize statement + proof.
\todoresearcher{Finalize the exact statement of Proposition T1
(masked-coalition value functions cannot distinguish certain
neighborhood substitutes without retraining) and write the proof.
Sketch: construct two parameterizations of the brief run realizing the
same
function with credit routed through either substitute; show
$v_{\mathrm{prune}}$ is identical on all coalitions for one
parameterization but not under retraining.}

\begin{proof}[Proof sketch (placeholder)]
TBD. % TODO(researcher)
\end{proof}

\subsection{Statement and Proof Sketch of Proposition T3 (Rank Floor)}
\label{app:proof-t3}

% TODO(researcher): finalize statement + proof.
\todoresearcher{Finalize the exact statement of Proposition T3 (the
rank floor) and write the full proof.}

\begin{theorybox}
\begin{propositionTrank}[Rank floor; statement to be finalized]
% (no \label: unnumbered T-statement; refer to it as "Proposition~T3")
Let $\mathcal{C}$ be a complex with $R$ ranks and $S$ a coalition of
neighborhoods. If the rank pairs connected by the neighborhoods in $S$
do not span the rank hierarchy, then every cell of each unspanned rank
has an empty composed receptive field: no input signal can reach it
through any sequence of message-passing layers, and each such cell is a
null player of the masked game. Spanning $R$ ranks requires on the
order of $R$ neighborhoods; hence $k \ge R$ is a structural floor on
the size of any coalition that uses the whole complex---$k = 3$ for the
2-complexes (nodes, edges, faces) of molecular benchmarks. \emph{Exact
statement TBD by researcher.}
\end{propositionTrank}
\end{theorybox}

\begin{proof}[Proof sketch]
Each neighborhood is an operator between one pair of ranks
(\cref{sec:prelim-gccn}), so a coalition $S$ induces a
rank-connectivity multigraph $G_S$ on the ranks of the complex, with
one edge per neighborhood in $S$. Composed receptive fields grow only
along walks of $G_S$: after any number of message-passing layers, a
cell of rank $r$ can have received signal only from cells whose ranks
lie in the connected component of $r$ in $G_S$. If a rank is
disconnected from every rank carrying input signal, all of its cells
have empty composed receptive fields; masking such a cell changes no
coalition value, so it is a null player and the null-player axiom
assigns it exactly zero. Connecting all ranks of an $R$-rank complex
requires a connected spanning subgraph of $G_S$, hence on the order of
$R$ neighborhood edges; accounting for the directionality of routes
and for which ranks carry input signal versus readout yields the floor
$k \ge R$ stated above and discussed in \cref{app:selection-rankfloor}
($k \ge 3$ for the 2-complexes of the molecular benchmarks).
% TODO(researcher): tighten the directionality/readout counting into
% the final constant.
\end{proof}

\subsection{Proof of Corollary T2 (Hypergraph Message Passing)}
\label{app:proof-t2}

% TODO(researcher): finalize statement + proof.
\todoresearcher{Finalize the exact statement of Corollary T2 (the
framework extends to hypergraph message passing; carried by theory only,
the hypergraph experiment was cut) and write the proof. Sketch:
hypergraph incidence defines the neighborhood set; players, masking, and
both value functions are unchanged; axioms follow from
\cref{eq:shapley,eq:interaction} verbatim.}

\begin{proof}[Proof sketch (placeholder)]
TBD. % TODO(researcher)
\end{proof}

% ======================================================================
% Appendix -- Architecture selection: details, analyses, and evidence
% structure. Depth moved out of the main text (researcher-set page
% budget); main-text sections reference these subsections.
% Binding style rules apply here too.
% ======================================================================
\section{Architecture Selection: Details and Evidence Structure}
\label{app:selection}

\subsection{Full Selection Results}
\label{app:selection-table}

\Cref{tab:performance-summary} tabulates every selection arm and sweep
marker of \cref{fig:lead,fig:lead-bench} with exact values and
per-row cost multiples.

% ======================================================================
% tables/performance_summary.tex
% *** AUTO-GENERATED by scripts/make_tables.py --performance-table ***
% *** from data/frozen/{autok_marker,ladder_b2_{marker,t1,mantra},hopse_k3_rung, ***
% *** baseline_markers,hopse_sweep_{t1,mantra}}. DO NOT EDIT BY HAND. ***
% ======================================================================
\begin{table}[t]
\centering
\caption{\textbf{Architecture selection.} Test accuracy
  (mean $\pm$ sd over seeds \unfrozen{\ResSeedList}) and, in
  parentheses, FLOPs as a multiple of the $B{=}1$ selection run
  (\cref{app:flops-conventions}). Bold marks the best
  accuracy in each row; gray marks entries within one sd of
  it. The last column quantifies the tradeoff: the accuracy
  of the best selection arm minus the best sweep's, and, in
  parentheses, that sweep's compute as a multiple of that
  selection arm's. Sweeps are priced at their full grid cost.
  --- = arm not defined or not run;
  $^{\dagger}$no $B{=}1$ arm, multiples relative to $B{=}2$;
  $^{e}$best of the 64-config retraining enumeration at
  estimated cost; $^{b}$balanced test accuracy (plain accuracy
  collapses to the majority class).}
\label{tab:performance-summary}
\small
\setlength{\tabcolsep}{3.5pt}
\resizebox{\linewidth}{!}{%
\begin{tabular}{lcccccc}
\toprule
 & \multicolumn{2}{c}{Ours (selection, GCCN)} & \multicolumn{1}{c}{Ours (selection, HOPSE)} & \multicolumn{2}{c}{Sweeps (at full grid cost)} & \multicolumn{1}{c}{Tradeoff} \\
\cmidrule(lr){2-3}\cmidrule(lr){4-4}\cmidrule(lr){5-6}\cmidrule(lr){7-7}
Dataset & $B{=}1$ & $B{=}2$ & $B{=}2$ & GCCN sweep & HOPSE sweep & $\Delta$ (sweep cost) \\
\midrule
\multicolumn{7}{l}{\emph{GraphXAI benchmarks (lifted to cellular complexes)}} \\
Benzene & \unfrozen{0.925}$\pm$0.007 {\scriptsize($\times$1.0)} & \unfrozen{0.926}$\pm$0.004 {\scriptsize($\times$2.0)} & \cellcolor{black!10}\unfrozen{0.934}$\pm$0.001 {\scriptsize($\times$5.4)} & \unfrozen{0.926}$\pm$0.002 {\scriptsize($\times$22.7)} & \cellcolor{black!10}\textbf{\unfrozen{0.936}$\pm$0.003} {\scriptsize($\times$36.4)} & \unfrozen{-0.002} {\scriptsize($\times$7)} \\
Fluoride-Carbonyl & \unfrozen{0.901}$\pm$0.057 {\scriptsize($\times$1.0)} & \cellcolor{black!10}\textbf{\unfrozen{0.939}$\pm$0.006} {\scriptsize($\times$1.9)} & \unfrozen{0.918}$\pm$0.004 {\scriptsize($\times$8.5)} & \cellcolor{black!10}\unfrozen{0.938}$\pm$0.004 {\scriptsize($\times$20.6)} & \cellcolor{black!10}\unfrozen{0.934}$\pm$0.009 {\scriptsize($\times$41.6)} & \unfrozen{+0.001} {\scriptsize($\times$11)} \\
Mutagenicity & \cellcolor{black!10}\unfrozen{0.967}$\pm$0.011 {\scriptsize($\times$1.0)} & \cellcolor{black!10}\unfrozen{0.975}$\pm$0.010 {\scriptsize($\times$2.0)} & \unfrozen{0.961}$\pm$0.008 {\scriptsize($\times$5.9)} & \unfrozen{0.964}$\pm$0.014 {\scriptsize($\times$15.4)} & \cellcolor{black!10}\textbf{\unfrozen{0.976}$\pm$0.009} {\scriptsize($\times$36.3)} & \unfrozen{-0.001} {\scriptsize($\times$18)} \\
Alkane-Carbonyl & \unfrozen{0.965}$\pm$0.006 {\scriptsize($\times$1.0)} & \unfrozen{0.965}$\pm$0.006 {\scriptsize($\times$2.2)} & \unfrozen{0.969}$\pm$0.009 {\scriptsize($\times$10.5)} & \cellcolor{black!10}\textbf{\unfrozen{0.984}$\pm$0.007} {\scriptsize($\times$95.0)}$^{e}$ & \unfrozen{0.963}$\pm$0.007 {\scriptsize($\times$47.8)} & \unfrozen{-0.016} {\scriptsize($\times$9)} \\
\midrule
\multicolumn{7}{l}{\emph{TopoBench benchmarks (lifted to cellular complexes)}} \\
NCI1 & \cellcolor{black!10}\textbf{\unfrozen{0.764}$\pm$0.011} {\scriptsize($\times$1.0)} & \cellcolor{black!10}\unfrozen{0.758}$\pm$0.010 {\scriptsize($\times$2.8)} & \cellcolor{black!10}\unfrozen{0.759}$\pm$0.024 {\scriptsize($\times$5.6)} & \cellcolor{black!10}\unfrozen{0.759}$\pm$0.022 {\scriptsize($\times$15.0)} & \unfrozen{0.745}$\pm$0.011 {\scriptsize($\times$34.5)} & \unfrozen{+0.005} {\scriptsize($\times$15)} \\
MUTAG & \unfrozen{0.879}$\pm$0.025 {\scriptsize($\times$1.0)} & \cellcolor{black!10}\unfrozen{0.915}$\pm$0.021 {\scriptsize($\times$2.0)} & \unfrozen{0.844}$\pm$0.012 {\scriptsize($\times$10.9)} & \cellcolor{black!10}\textbf{\unfrozen{0.922}$\pm$0.033} {\scriptsize($\times$124.5)}$^{e}$ & \unfrozen{0.844}$\pm$0.054 {\scriptsize($\times$53.3)} & \unfrozen{-0.007} {\scriptsize($\times$61)} \\
NCI109 & \cellcolor{black!10}\unfrozen{0.758}$\pm$0.011 {\scriptsize($\times$1.0)} & \cellcolor{black!10}\unfrozen{0.755}$\pm$0.011 {\scriptsize($\times$1.7)} & \unfrozen{0.720}$\pm$0.005 {\scriptsize($\times$5.2)} & \cellcolor{black!10}\textbf{\unfrozen{0.762}$\pm$0.007} {\scriptsize($\times$38.8)}$^{e}$ & \unfrozen{0.737}$\pm$0.011 {\scriptsize($\times$27.1)} & \unfrozen{-0.004} {\scriptsize($\times$39)} \\
Cora & \unfrozen{0.838}$\pm$0.008 {\scriptsize($\times$1.0)} & \unfrozen{0.835}$\pm$0.005 {\scriptsize($\times$2.3)} & --- & \cellcolor{black!10}\textbf{\unfrozen{0.856}$\pm$0.006} {\scriptsize($\times$14.9)} & --- & \unfrozen{-0.019} {\scriptsize($\times$15)} \\
Citeseer & \unfrozen{0.673}$\pm$0.014 {\scriptsize($\times$1.0)} & \unfrozen{0.670}$\pm$0.007 {\scriptsize($\times$1.8)} & --- & \cellcolor{black!10}\textbf{\unfrozen{0.683}$\pm$0.002} {\scriptsize($\times$15.1)} & --- & \unfrozen{-0.010} {\scriptsize($\times$15)} \\
Pubmed & \unfrozen{0.860}$\pm$0.004 {\scriptsize($\times$1.0)} & \unfrozen{0.863}$\pm$0.004 {\scriptsize($\times$2.3)} & --- & \cellcolor{black!10}\textbf{\unfrozen{0.865}$\pm$0.001} {\scriptsize($\times$14.4)} & --- & \unfrozen{-0.002} {\scriptsize($\times$6)} \\
\midrule
\multicolumn{7}{l}{\emph{Natively higher-order (simplicial)}} \\
MANTRA$^{b}$ & \unfrozen{0.532}$\pm$0.028 {\scriptsize($\times$1.0)} & \unfrozen{0.615}$\pm$0.103 {\scriptsize($\times$1.6)} & \unfrozen{0.771}$\pm$0.017 {\scriptsize($\times$9.6)} & \unfrozen{0.548}$\pm$0.000 {\scriptsize($\times$11.0)} & \cellcolor{black!10}\textbf{\unfrozen{0.792}$\pm$0.009} {\scriptsize($\times$68.7)} & \unfrozen{-0.021} {\scriptsize($\times$7)} \\
\bottomrule
\end{tabular}}
\end{table}


\subsection{The Candidate Sequence, Formally}
\label{app:selection-ladder}

The players are the family's neighborhood-indexed components (GCCN
message-passing routes; HOPSE encoding blocks;
\cref{sec:framework-game2}). Backward elimination starts from the full
coalition and repeatedly removes the component whose removal is
cheapest on $v_{\mathrm{prune}}$, while that cost is below the
validation binomial floor $\sqrt{p(1-p)/m}$. Because each architecture
is nested in the previous one, the eliminations yield one candidate
per size. The candidate choice was fixed before any $B \ge 2$
experiment ran. The \emph{stop candidate} is the smallest
architecture whose next removal would cost more than the floor.
Candidates $2..B$ are the architectures with the highest cheap-game
value at sizes not yet taken. $B{=}1$ trains the stop candidate only
($\sim$\unfrozen{\ResLadderCostBOne} training-run equivalents);
$B{=}2$ adds one candidate ($\sim$\unfrozen{\ResLadderCostBTwo});
$B{=}\unfrozen{\ResMenuEntries}$ with literature architectures
recovers the menu; $B{=}\text{all}$ is the sweep. Exact protocol in
\cref{app:pipeline-protocol}.

\subsection{Why Not a Sorted-Shapley Elbow: Substitute-Commitment}
\label{app:selection-commitment}

Sorting components by Shapley value and cutting at the elbow fails
under \emph{redundancy}: when two components can do the same job, the
briefly trained full model commits to one (which one varies with the
training seed), and the attribution faithfully reports that model's
realized division of labor rather than what a retrained model could
do (Proposition~T1). We
call this \emph{substitute-commitment}; it independently reproduces
the valuation-is-not-selection result of \citet{fryer2021shapley} in
the wild. Pooled correlation between the brief run's Shapley reliance
and post-pruning viability is weak: $\rho =$ \unfrozen{\ResPooledRho}
(95\% CI [\unfrozen{\ResPooledRhoCILow},
\unfrozen{\ResPooledRhoCIHigh}]), consistent with redundancy rather
than a broken estimator (exact to \unfrozen{\ResExactnessMaxErr} on
its game). The honest elbow is the greedy \emph{trace}: the sequence
of measured removal costs along the elimination path, in which a
redundant component is
removed for free while the last copy of a function visibly hurts and
stays.

\begin{figure}[h]
  \centering
  \IfFileExists{../figures/out/substitute_commitment.pdf}{%
    \includegraphics[width=\linewidth]{../figures/out/substitute_commitment.pdf}%
  }{%
    \figureplaceholder{figures/out/substitute_commitment.pdf}%
  }
  \caption{Substitute-commitment on the Benzene exemplar. Left:
    per-seed reliance of the briefly trained full model on
    \nbhd{up\_adjacency-0} vs.\ its incidence substitutes; one seed
    commits to the adjacency route
    while the other two lean on an incidence substitute; attribution
    follows the commitment. Right: post-pruning viability on the exact
    landscape (leave-one-out and best-subset-without values against
    the full-vocabulary band): removing the committed route costs at
    most $-0.0033$ test accuracy ($\approx$1.3 seed-sd), so the
    commitment is not a necessity. Subset markers carry the record's
    median seed-sd as a representative whisker. Linear axes.}
  \label{fig:substitute-commitment}
\end{figure}

\subsection{Why Backward, Not Forward: The Plateau}
\label{app:selection-backward}

Growing an architecture from empty fails on briefly trained models: a
\unfrozen{\ResStemEpochs}-epoch full model masked down to one or
two components behaves like an untrained model: small-coalition
values collapse to one-class accuracy (the score of predicting a
single class for everything), marginal gains read zero, and forward
greedy stops on arbitrary singletons (Fluoride-Carbonyl
\unfrozen{\ResForwardFluoride} vs.\ the
\unfrozen{\ResForwardFluorideRef} reference; Cora
\unfrozen{\ResForwardCoraMean} $\pm$ \unfrozen{\ResForwardCoraSd};
exploratory development runs, labeled as such). The plateau replicated
independently on two clusters and in a second model family
(\cref{app:trust-details}). The rule itself is sound; on the
\emph{true} retraining game the same greedy stop is the best rule we
scored, so the walk, not the rule, was at fault. Backward elimination
starts at the full model, where the brief run's game is faithful, and only
interrogates that region: it beats forward greedy on
\unfrozen{\ResBackwardBeatsForwardWins}/\unfrozen{\ResBackwardBeatsForwardTotal}
datasets and collapses its seed variance (Cora sd
\unfrozen{\ResForwardCoraSd} $\to$ \unfrozen{\ResBackwardCoraSd};
exploratory).

\subsection{The Rank Floor: Discussion}
\label{app:selection-rankfloor}

Proposition~T3 (statement and proof sketch in \cref{app:proof-t3})
explains the literature: HOPSE and CWN shipped defaults are exactly
rank-spanning 3-sets, so fixed $k=3$ looks strong on 2-complex
benchmark suites: the floor plus a small surplus. The dark blue star
of \cref{fig:lead-full} is our fixed $k=3$ variant at its true one-run
cost, strong on these suites for exactly this reason. Off-suite the
floor is not a law: on MUTAG the true optimal size is
\unfrozen{\ResMutagTrueKStar} and fixed $k=3$ loses
\unfrozen{\ResFixedKMutagRegret} against the exact landscape optimum
(mean regret \unfrozen{\ResFixedKMeanRegret} across the
\unfrozen{\ResLadderAllTotal} exact landscapes)---the caution against
reading ``$k=3$ works'' as universal. The campaign measured that tail
live: the anchored $k{=}3$ arm lands at
\unfrozen{\ResMutagFixedKLiveAcc} $\pm$
\unfrozen{\ResMutagFixedKLiveSd} on MUTAG while the $B{=}2$ selection
holds \unfrozen{\ResMutagLadderBTwoAcc} $\pm$
\unfrozen{\ResMutagLadderBTwoSd} against the exact optimum
\unfrozen{\ResMutagExactOpt} (\cref{sec:performance}). The selection
recovers the floor where it is optimal and grows past it where it is
not.

\subsection{Evidence Structure: Frozen, Exploratory, Pre-Registered}
\label{app:selection-evidence}

Selection rules are easy to overfit to the datasets used while
developing them; we keep three evidence layers separate.
\textbf{Frozen rule}: the elimination path and its candidate choice
were fixed before
any $B \ge 2$ experiment ran and before the out-of-sample scoring; no
knob was revisited afterwards. \textbf{Out-of-sample validation}:
\unfrozen{\ResLadderOOSTotal} exhaustively-enumerated architecture
landscapes (MUTAG, NCI109, ZINC, ZINC-GIN, alkane-carbonyl) were never
used during development; the candidate sequence contains the exact
global optimum on \unfrozen{\ResLadderOOSContained}/\unfrozen{\ResLadderOOSTotal}
(across all \unfrozen{\ResLadderAllTotal} landscapes,
\unfrozen{\ResLadderAllContained}/\unfrozen{\ResLadderAllTotal}
exact-zero regret; worst \unfrozen{\ResLadderOOSWorstRegret}),
validating the \emph{shape} of the rule on true retraining games.
\textbf{Exploratory vs.\ pre-registered}: the
\unfrozen{\ResMarkerDatasets}-dataset development suite guided the
design and its numbers are labeled exploratory; masked-game quality on
unseen data is the job of the pre-registered confirmatory battery
(protocol in \cref{app:experimental}), whose verdicts are reported
as-is. Development heuristics explored and retired are disclosed in
\cref{app:dev-trajectory}.

\subsection{Baselines: Three Layers of ``Defaults''}
\label{app:selection-baselines}

``The default'' is not one thing in TDL; the baseline field has three
layers, and \cref{fig:lead-full} together with
\cref{tab:performance-summary} represents all three. \textbf{Hardwired architecture constants}: CWN
\citep{bodnar2021cwn}, CCXN \citep{hajij2020ccxn}, CAN
\citep{giusti2022can}, SCN and SCCN \citep{yang2022sccnn} fix their
neighborhood sets in the architecture definition. \textbf{Shipped
configuration defaults}: TopoTune's
\unfrozen{\ResTopoTuneRouteDefault}-route default
\citep{papillon2025topotune}; HOPSE's
\unfrozen{\ResHopseModelDefaultBlocks}-encoding model default vs.\ its
\unfrozen{\ResHopseExptDefaultBlocks}-encoding experiment default
\citep{hopse2025}. The \emph{menu} marker is this layer: the
\unfrozen{\ResMenuEntries} published fixed architectures priced
together at their true cost, selected per-seed by validation.
\textbf{Paper-time sweeps}: TopoTune's
\unfrozen{\ResTopoTuneSweepStructures} neighborhood structures crossed
with sub-models; HOPSE's
\unfrozen{\ResHopseSweepConfigs}-configuration sweep---the pale
markers, each at the cost of running its whole grid once. Cost
accounting is uniform across all of them
(\cref{app:flops-conventions}); HOPSE's preprocessing is charged to
HOPSE and is FLOPs-cheap but wall-clock-heavy, so the FLOPs axis is
the one generous to the baselines (wall-clock variant in
\cref{app:figures}).

% ======================================================================
% Appendix -- Experimental details
% Binding style rules apply here too: multi-seed error bars,
% parameter/budget matching on all comparisons, "neighborhoods".
% ======================================================================
\section{Experimental Details}
\label{app:experimental}

\subsection{Datasets and Liftings}
% TODO(researcher): dataset statistics table (graphs, cells per rank
% after lifting, class balance) for Benzene, Fluoride-Carbonyl,
% Mutagenicity, MANTRA, MUTAG, NCI1, citation graphs.
\todoresearcher{Fill dataset statistics and lifting configuration
(TopoBench lifting names + parameters).}

\subsection{Dataset Coverage and Scope}
\label{app:coverage}

\textbf{Coverage.} The selection experiments cover
\unfrozen{\ResLeadPanels} datasets: the
\unfrozen{\ResMarkerDatasets}-dataset development suite (Benzene,
Fluoride-Carbonyl, Mutagenicity, NCI1, Cora, Citeseer, Pubmed) plus
four datasets added at campaign time (MUTAG, NCI109, Alkane-Carbonyl,
MANTRA orientability), every arm under the one pipeline of
\cref{app:pipeline-protocol} at the shared seeds
\unfrozen{\ResSeedList} with measured FLOPs. Every row of
\cref{tab:performance-summary} carries a $B{=}1$ arm; runs that
predate the profiler are priced post hoc from their recorded
architectures (\cref{app:flops-conventions}). Where an
exhaustively-enumerated exact landscape exists (MUTAG, NCI109,
Alkane-Carbonyl; \cref{app:selection-evidence}) its optimum stands in
for the GCCN sweep, marked $^{e}$ in the table.

\textbf{Out of scope, declared.} No selection arm ran on PROTEINS;
it appears in the faithfulness study only (\cref{tab:nogt}). The
HOPSE family is evaluated on complex-level tasks; the node-level
citation complexes carry the GCCN family only. HOPSE's clinical quartet (BBB-Martins, Caco2-Wang,
CYP3A4-Veith, Clearance-Hepatocyte-AZ) is likewise out of scope: none
of our arms ran there, and we report no borrowed numbers. ZINC is
\emph{deferred}, not skipped: the stop rule derives from a
classification validation set's binomial noise floor, ZINC is a
regression task, and its arm waits on a pre-registered regression
noise floor rather than an analogue improvised after seeing results
(its exact landscape already serves in the out-of-sample scoring of
\cref{app:selection-evidence}).

\textbf{The MANTRA metric.} MANTRA orientability is class-imbalanced:
the majority class is \unfrozen{\ResMantraMajorityRate} of every
split (per-seed class-balance record frozen with the results). An
executable sanity check comparing each cell's plain accuracy against
that rate flagged \unfrozen{\ResMantraCollapsedModels} frozen cells
as byte-identical majority collapses, so every MANTRA number in the
paper is balanced test accuracy (mean per-class recall), re-scored
from each arm's saved best checkpoint. Selection inside every arm
still uses plain validation accuracy exactly as wired---the balanced
metric never re-selects---so the re-scoring changes no pick, only the
reported score.

\subsection{Models and Training}
\label{app:models-training}

\textbf{Model configuration (architecture selection, GCCN family).} All architecture-%
selection experiments use the \texttt{topotune\_nonlinear\_exact}
configuration: the TopoTune GCCN backbone with a linear sum-pool readout
(\texttt{NoReadOut}). \textbf{Protocol-fidelity calibration}: at the
Benzene calibration architecture (mask \unfrozen{\ResCalibMaskID}), the
backbone parameter count reproduces the recorded experimental protocol
exactly---\unfrozen{\ResCalibBackboneParams} backbone parameters---which
is how we verify that reimplemented runs train the same model as the
original record, down to the readout choice.

\textbf{Seeds.} Every comparison uses the identical seed set
\unfrozen{\ResSeedList} for our pipeline and for every baseline; on
the seven development-suite datasets our own arms additionally carry
$n=$ \unfrozen{\ResOurReplications} seeds in the frozen record.
Selection within every method is per-seed
validation-best; no method sees test data before final scoring.
% TODO(researcher): optimizer, LR schedule, epochs, early stopping,
% hardware, wall-clock per run.
\todoresearcher{Fill remaining training configuration; must match the
budget-matching protocol referenced in \cref{sec:predictions}.}

\subsection{Protocol Uniformity Across Methods}
\label{app:protocol-uniformity}

One selection pipeline is applied identically across every dataset and
every model family; the table below is generated directly from the run
manifests so that uniformity is checkable rather than asserted.
% TODO(generated): scripts/make_tables.py will emit
% tables/protocol_uniformity.tex from data/frozen/; the current file is
% a placeholder stub.
% ======================================================================
% tables/protocol_uniformity.tex
% *** TODO(generated): scripts/make_tables.py will emit this table   ***
% *** from data/frozen/ run manifests. The current contents are a    ***
% *** hand-written placeholder stub; do not hand-fill values.        ***
% ======================================================================
\begin{table}[h]
\centering
\caption{Protocol uniformity across methods and families: one selection
  pipeline, one seed set, one cost convention.
  \todoresearcherinline{TODO(generated): emitted by
  scripts/make\_tables.py from run manifests.}}
\label{tab:protocol-uniformity}
\small
\begin{tabular}{lll}
\toprule
Protocol element & Ours (GCCN / HOPSE) & Baselines \\
\midrule
Seeds & \unfrozen{\ResSeedList} & \unfrozen{\ResSeedList} \\
Selection rule & per-seed validation-best & per-seed validation-best \\
FLOPs convention & \multicolumn{2}{l}{identical (\cref{app:flops-conventions})} \\
Brief-run epochs & \unfrozen{\ResStemEpochs} & --- \\
Attribution budget & \unfrozen{\ResSampledPasses} passes
  ($\sim$\unfrozen{\ResUniqueEvals} unique evals) & --- \\
\ResTBD & \ResTBD & \ResTBD \\
\bottomrule
\end{tabular}
\end{table}


\subsection{FLOPs Measurement Conventions}
\label{app:flops-conventions}

All FLOPs in this paper are measured under one convention, identical
across our pipeline and every baseline:
\begin{itemize}
  \item PyTorch \texttt{FlopCounterMode} with dense pricing;
  \item explicit shims for operations the counter does not price
    natively: eigendecomposition (\texttt{eigh}), SVD, and sparse
    operations;
  \item every stage of every method is charged, including HOPSE's
    encoding preprocessing. That preprocessing is FLOPs-cheap but
    wall-clock-heavy, so the FLOPs axis of \cref{fig:lead} is the
    choice \emph{generous to the baselines}; the wall-clock variant of
    the figure is in \cref{app:figures}.
\end{itemize}
Cost claims in training-run equivalents are priced on unique coalition
evaluations (\cref{app:accounting}); the per-dataset,
per-stage ledger is generated from the frozen measurements.

Selection runs that predate the profiler are priced post hoc: we
rebuild the recorded architectures, profile one training epoch and the
evaluation passes under the convention above, and combine the counts
with the recorded epoch and evaluation tallies. FLOP counts depend
only on architecture and data, so these costs are measured, not
estimated.
% TODO(researcher): frozen cost-ledger table pointer once generated.

\subsection{The Canonical Pipeline: Exact Protocol}
\label{app:pipeline-protocol}

Brief run: the all-components model
(\unfrozen{\ResNumNeighborhoods} neighborhoods in the GCCN family; all
encoding blocks in HOPSE), trained \unfrozen{\ResStemEpochs} epochs.
Cheap game: masked validation accuracy with removal semantics matching
deployment-time pruning (\cref{sec:framework-game2}); attribution
budget \unfrozen{\ResSampledPasses} passes
($\sim$\unfrozen{\ResUniqueEvals} unique coalition evaluations after
caching). Elimination: backward from the full coalition,
removing the cheapest component while removal cost is below the
1-standard-error binomial floor $\sqrt{p(1-p)/m}$ of the validation
set. Candidates: candidate 1 = the stop candidate; candidates $2..B$
= highest-cheap-value coalitions at sizes not yet taken (rule frozen
before any $B \ge 2$ experiment ran). Continuation: train each
selected candidate to the standard budget; select by validation
accuracy.
% TODO(researcher): exact epoch counts, sampling seeds, and the
% measured ledger backing the \ResLadderCostBOne / \ResLadderCostBTwo
% training-run-equivalent claims.

\subsection{Development Trajectory (Disclosure)}
\label{app:dev-trajectory}

For evaluation hygiene we disclose the selection heuristics explored
before the rule was frozen. During development on the
\unfrozen{\ResMarkerDatasets}-dataset suite we iterated on (i)
\emph{anchored selection}, which always keeps the
\nbhd{up\_adjacency-0} neighborhood---effective on the molecule
datasets but contradicted on Fluoride-Carbonyl, where anchor-free
architectures win without any catastrophe, so the anchor is a
suite-specific prior rather than a rule; and (ii) a
\emph{certify-or-repair} fallback wrapper, which retrains a second
candidate when the first falls below an in-run reference. Both were
retired as fitted to the development suite. The frozen procedure
subsumes their useful behavior structurally: backward elimination
keeps last-copy components without being told which ones are anchors,
and budget $B=2$ provides the fallback as an explicit, pre-priced
candidate rather than a triggered repair. All numbers from the development suite
are labeled exploratory wherever they appear; the pre-registered
battery is the confirmatory layer (\cref{app:selection-evidence}).

\subsection{Baseline Explainers and Budget Matching (Prediction Explanations)}
% TODO(researcher): SubgraphX rollout budget, GNNExplainer /
% PGExplainer epochs and hyperparameters, and the shared evaluation
% budget that makes comparisons parameter/budget matched.
\todoresearcher{Fill prediction-explanation baseline budgets.}

\subsection{Trustworthiness Details}
\label{app:trust-details}

\textbf{What the axiom checks caught} (each independently confirmed and
reported upstream):
\begin{itemize}
  \item \textbf{Dead inter-rank routes}: swapped \texttt{edge\_index}
    rows silently zeroed messages along certain inter-rank
    neighborhoods; flagged by null-player checks assigning exact zeros
    to neighborhoods that should have mattered.
    % TODO(researcher): upstream issue/PR numbers.
  \item \textbf{Batch-dependent collapsed backbone}: a backbone
    configuration collapsed to batch-size-dependent behavior; flagged
    by attribution non-invariance across batch compositions.
  \item \textbf{Rank-aliasing offsets}: cell-index offsets aliased
    features across ranks; flagged by symmetry violations between
    provably symmetric cells.
  \item \textbf{Composed cross-rank batching bug}: individually correct
    components composed into incorrect cross-rank batching; caught by
    the invariance guard that re-runs attributions under permuted
    batching and requires bitwise-stable results.
\end{itemize}
% TODO(researcher): frozen check-suite runtime numbers + pointer to the
% released test suite.

\textbf{Second-family replication of substitute-commitment.} On
Mutagenicity (seed 43) the briefly trained HOPSE model commits to a
different trio of
encoding blocks than its sibling seeds; the committed pick lands
\unfrozen{\ResHopseCommitDropPts} accuracy points below them, and
validation already shows it (\unfrozen{\ResHopseCommitValBad} vs.\
\unfrozen{\ResHopseCommitValSibA} and
\unfrozen{\ResHopseCommitValSibB}). The mechanism is a property of
briefly trained models under redundancy, not of message passing in
particular.

\textbf{Catastrophe audit (validation visibility).} Across the
\unfrozen{\ResCatastropheCells} development-grid runs, exactly
\unfrozen{\ResCatastropheCount} selections were catastrophic (test
accuracy far below their replicate median); both were equally visible
in validation (\unfrozen{\ResCatastropheValA} and
\unfrozen{\ResCatastropheValB} against
$\sim$\unfrozen{\ResCatastropheSiblingVal} for siblings), with
\unfrozen{\ResCatastropheFalseAlarms} false alarms at the same
threshold---which is why the pipeline's final decision is always a
validation comparison of trained candidates.

\textbf{The plateau as an executable check.} On
\unfrozen{\ResStemEpochs}-epoch briefly trained models the cheap
game's gain
\emph{magnitudes} saturate---in the cross-cluster replication,
\unfrozen{\ResPlateauZeroCells}/\unfrozen{\ResPlateauTotalCells} runs
had exactly-zero best gains at small coalitions---while the
\emph{ordering} survives. Exact zeros in gain magnitudes are an
executable signature, not a judgment call, which is how the
forward-greedy null (\cref{app:selection-backward}) was diagnosed
rather than merely observed; the design consequence is that every
threshold decision consumes trained-candidate validation or the
binomial floor, never a raw cheap-game gain.

\subsection{Auditable Cost Accounting: Unique Coalition Evaluations}
\label{app:accounting}

The ``\unfrozen{\ResSampledPasses}-pass'' attribution budget counts
\emph{unique} coalition evaluations:
$\sim$\unfrozen{\ResUniqueEvals} cached evaluations, or
$\approx$\unfrozen{\ResUniqueEvalFraction}$\times$ the exact game's
\unfrozen{\ResExactGameCoalitions} coalitions. Permutation sampling
revisits coalitions; with caching, the number of forward passes
actually paid is the unique-evaluation count, not the nominal pass
count---the two can differ by an order of magnitude---and the measured
wall clocks confirm the accounting, so every cost claim in
\cref{sec:performance} is priced on the paid number.

\subsection{Reproducibility}
\label{app:reproducibility}

Every number in this paper is a macro regenerated from a frozen,
checksummed results directory; no result is typed by hand. One command
(\texttt{scripts/reproduce.sh}) reruns the test suite, regenerates all
tables, macros, and figures, and verifies presentation numbers against
the same frozen files. The freeze manifest hash for this draft is
\texttt{\ResManifestHash}; any post-freeze drift in the underlying data
fails the build.

\subsection{Axiom Check Suite}
% TODO(researcher): list each executable check (efficiency, null
% player, symmetry, batching/ordering invariance guard), its tolerance,
% and the reproduction commands for the
% \ResNumBugsTotal caught bugs (upstream issue links).
\todoresearcher{Document the executable-check suite and bug
reproductions.}

% ======================================================================
% Appendix C -- Per-dataset tables
% All tables in this appendix are AUTO-GENERATED by
% scripts/make_tables.py from data/frozen/ -- do not edit by hand.
% ======================================================================
\section{Per-Dataset Tables}
\label{app:tables}

% TODO(researcher): frozen data must land for every table below.

\subsection{Explanation Quality (Prediction Explanations)}
\label{app:tables-by-subject}

% AUTO-GENERATED by scripts/make_paper_tables.py -- do not edit
\begin{table}[t]
\centering
\small
\caption{\textbf{Per-subject explanation accuracy.} Node-level GEA (mean $\pm$ sd, 3 seeds) for each parameter-matched GNN subject; \cref{tab:graph-benchmarks} reports each method's better subject. The zero-fill row runs our explanation search with feature zero-filling, SubgraphX's removal semantics, in place of detachment. Bold marks the best entry per column; gray marks entries within one sd of it.}
\label{tab:by-subject}
\setlength{\tabcolsep}{3.5pt}
\resizebox{\linewidth}{!}{%
\begin{tabular}{lcccccccc}
\toprule
 & \multicolumn{2}{c}{Benzene} & \multicolumn{2}{c}{Fluoride} & \multicolumn{2}{c}{Mutagenicity} & \multicolumn{2}{c}{Alkane} \\
\cmidrule(lr){2-3}\cmidrule(lr){4-5}\cmidrule(lr){6-7}\cmidrule(lr){8-9}
Method & GCN & GIN & GCN & GIN & GCN & GIN & GCN & GIN \\
\midrule
\method{} (attribution) & 0.467 $\pm$ 0.078 & 0.859 $\pm$ 0.009 & \cellcolor{black!10}0.398 $\pm$ 0.040 & \cellcolor{black!10}0.507 $\pm$ 0.047 & \cellcolor{black!10}\textbf{0.723 $\pm$ 0.074} & \cellcolor{black!10}\textbf{0.798 $\pm$ 0.124} & \cellcolor{black!10}\textbf{0.411 $\pm$ 0.058} & \cellcolor{black!10}\textbf{0.500 $\pm$ 0.081} \\
\method{} (explanation) & 0.378 $\pm$ 0.012 & \cellcolor{black!10}\textbf{0.969 $\pm$ 0.009} & 0.290 $\pm$ 0.095 & 0.340 $\pm$ 0.011 & 0.319 $\pm$ 0.067 & 0.480 $\pm$ 0.135 & 0.318 $\pm$ 0.073 & 0.354 $\pm$ 0.125 \\
\method{} (explanation, zero-fill) & \cellcolor{black!10}0.648 $\pm$ 0.055 & 0.569 $\pm$ 0.064 & 0.244 $\pm$ 0.034 & 0.313 $\pm$ 0.038 & 0.561 $\pm$ 0.062 & 0.273 $\pm$ 0.164 & \cellcolor{black!10}0.378 $\pm$ 0.075 & 0.391 $\pm$ 0.028 \\
SubgraphX & \cellcolor{black!10}\textbf{0.664 $\pm$ 0.027} & 0.796 $\pm$ 0.029 & 0.226 $\pm$ 0.024 & 0.439 $\pm$ 0.107 & 0.344 $\pm$ 0.036 & 0.274 $\pm$ 0.055 & 0.276 $\pm$ 0.013 & 0.296 $\pm$ 0.010 \\
PGExplainer & \cellcolor{black!10}0.642 $\pm$ 0.044 & 0.453 $\pm$ 0.239 & \cellcolor{black!10}\textbf{0.429 $\pm$ 0.092} & \cellcolor{black!10}\textbf{0.555 $\pm$ 0.067} & 0.290 $\pm$ 0.102 & 0.175 $\pm$ 0.200 & 0.203 $\pm$ 0.137 & 0.205 $\pm$ 0.160 \\
GNNExplainer & 0.443 $\pm$ 0.054 & 0.256 $\pm$ 0.081 & 0.210 $\pm$ 0.014 & 0.166 $\pm$ 0.035 & 0.152 $\pm$ 0.033 & 0.158 $\pm$ 0.023 & 0.066 $\pm$ 0.050 & 0.162 $\pm$ 0.049 \\
Random & 0.250 $\pm$ 0.000 & 0.250 $\pm$ 0.000 & 0.219 $\pm$ 0.000 & 0.219 $\pm$ 0.000 & 0.132 $\pm$ 0.000 & 0.132 $\pm$ 0.000 & 0.173 $\pm$ 0.000 & 0.173 $\pm$ 0.000 \\
\bottomrule
\end{tabular}}
\end{table}


% \input{tables/gef_per_dataset}          % TODO: generate
% \input{tables/groundtruth_agreement}    % TODO: generate
% \input{tables/faithfulness_controls}    % TODO: generate
\todoresearcher{Generate per-dataset GEF, ground-truth agreement, and
controls tables via scripts/make\_tables.py.}

\subsection{Architecture Selection}
% \input{tables/selection_per_dataset}    % TODO: generate (full 7-dataset version, GCCN + HOPSE arms)
% \input{tables/candidates_per_dataset} % TODO: generate (candidate sizes, cheap values, val/test per candidate)
% \input{tables/ours_replications}        % TODO: generate (n=5-seed replications of our arms)
% \input{tables/exact_vs_sampled}         % TODO: generate (rho, regret, cost)
% \input{tables/rule_scoreboard}          % TODO: generate (auto-k rules vs exact truth, incl. binomial vs oracle stop)
% \input{tables/unique_eval_ledger}       % TODO: generate (unique coalition evaluations + wall clocks per stage)
% \input{tables/cost_ledger}              % TODO: generate (measured training-run equivalents)
\todoresearcher{Generate architecture-selection per-dataset tables via
scripts/make\_tables.py.}

\subsection{MANTRA Diagnosis}
% \input{tables/mantra_surrogate}         % TODO: generate (per-task surrogate agreement)
% \input{tables/mantra_hodge}             % TODO: generate (Hodge-feature accuracy vs count-only score)
\todoresearcher{Generate MANTRA surrogate and Hodge-feature tables.}


% ======================================================================
% Appendix D -- Additional figures
% Binding style rules: linear axes; uniform y-span across panels; error
% bars on every marker; one star per fixed k; menu = single marker at
% its true cost.
% ======================================================================
\section{Additional Figures}
\label{app:figures}


\begin{figure}[h]
  \centering
  \IfFileExists{../figures/out/lead_figure_full.pdf}{%
    \includegraphics[width=\linewidth]{../figures/out/lead_figure_full.pdf}%
  }{%
    \figureplaceholder{figures/out/lead_figure_full.pdf}%
  }
  \caption{Complete-marker variant of the lead figure:
    identical data and protocol as \cref{fig:lead}, adding the markers
    hidden there for readability --- the
    \unfrozen{\ResMenuEntries}-entry menu of published fixed
    architectures at its true cost (six
    trainings), the fixed $k{=}6$ variant of our pipeline, and the
    budget-$1$ arm. Every marker follows the same conventions
    (multi-seed error bars, per-seed validation-best selection, true
    costs).}
  \label{fig:lead-full}
\end{figure}

\begin{figure}[h]
  \centering
  \IfFileExists{../figures/out/lead_figure_wall.pdf}{%
    \includegraphics[width=\linewidth]{../figures/out/lead_figure_wall.pdf}%
  }{%
    \figureplaceholder{figures/out/lead_figure_wall.pdf}%
  }
  \caption{Wall-clock variant of the lead figure: same
    markers and protocol as \cref{fig:lead}, horizontal axis measured
    wall-clock seconds instead of FLOPs. HOPSE preprocessing is
    wall-clock-heavy, so this is the axis less generous to the
    baselines (\cref{app:flops-conventions}).}
  \label{fig:app-lead-wall}
\end{figure}

\begin{figure}[h]
  \centering
  \IfFileExists{../figures/out/estimator_variance.pdf}{%
    \includegraphics[width=\linewidth]{../figures/out/estimator_variance.pdf}%
  }{%
    \figureplaceholder{figures/out/estimator_variance.pdf}%
  }
  \caption{Selection-level cost of the sampled game vs.\ number of
    permutation passes: cheap-value regret of the sampled game's picks
    against the exact game's selection (the quantity behind the
    committed \unfrozen{\ResRegretMax} bound), for budgets
    32--512, $k \in \{2, 3, 6\}$, mean $\pm$ sd over seeds
    \unfrozen{\ResSeedList}. Dashed line: the production 128-pass
    budget. Linear axes; uniform y-span across panels.}
  \label{fig:app-estimator-variance}
\end{figure}

\begin{figure}[h]
  \centering
  \IfFileExists{../figures/out/exhaustive_landscape.pdf}{%
    \includegraphics[width=\linewidth]{../figures/out/exhaustive_landscape.pdf}%
  }{%
    \figureplaceholder{figures/out/exhaustive_landscape.pdf}%
  }
  \caption{Exhaustive neighborhood-subset landscapes (12 datasets);
    ground truth for the stop-rule scoring of
    \cref{sec:performance} and the out-of-sample selection validation of
    \cref{app:selection-evidence}. Gray points: all 64 enumerated
    subsets at their coalition-mean value vs.\ subset size; line: the
    best-per-$k$ oracle; diamond: the full vocabulary with its
    across-seed sd (subset markers carry the record's median seed-sd
    as a representative whisker; per-seed per-coalition values live in
    the experiment record). PROTEINS has no enumerated landscape
    (excluded a priori). Linear axes; uniform y-span across panels.}
  \label{fig:app-landscape}
\end{figure}

\begin{figure}[h]
  \centering
  \IfFileExists{../figures/out/chemistry_case_studies.pdf}{%
    \includegraphics[width=\linewidth]{../figures/out/chemistry_case_studies.pdf}%
  }{%
    \figureplaceholder{figures/out/chemistry_case_studies.pdf}%
  }
  \caption{Cell-native attributions on four molecular case studies
    (atoms, bonds, and ring cells colored by Shapley attribution;
    ground-truth motif atoms outlined). Benzene's ring cells carry
    concentrated rank-2 credit; on Fluoride-Carbonyl the carbonyl bond
    dominates while the fluorine site stays near zero. The
    Mutagenicity panel is drawn without element labels (the frozen
    record ships none for that dataset).}
  \label{fig:chemistry-cases}
\end{figure}

\begin{figure}[h]
  \centering
  \IfFileExists{../figures/out/seed_traces.pdf}{%
    \includegraphics[width=\linewidth]{../figures/out/seed_traces.pdf}%
  }{%
    \figureplaceholder{figures/out/seed_traces.pdf}%
  }
  \caption{Per-seed exact retraining-game Shapley traces across the
    six neighborhoods (seeds \unfrozen{\ResSeedList}; open circles),
    with mean $\pm$ SEM markers, extending
    \cref{fig:substitute-commitment} to all 12 landscape datasets;
    the adjacency/incidence substitute trio is shaded. Linear axes;
    uniform y-span across panels.}
  \label{fig:app-seed-traces}
\end{figure}

\begin{figure}[h]
  \centering
  \IfFileExists{../figures/out/mantra_attributions.pdf}{%
    \includegraphics[width=\linewidth]{../figures/out/mantra_attributions.pdf}%
  }{%
    \figureplaceholder{figures/out/mantra_attributions.pdf}%
  }
  \caption{MANTRA attribution values on a matched test pair (a torus
    and a Klein-bottle-family surface with identical cell counts
    9/27/18): per-cell contribution to the orientability score for the
    three measured models. The menu model spreads identical negative
    credit over vertices on both complexes; the selected GCCN
    routes small credit through triangles only; HOPSE (own y-scale,
    stated on the axis) separates the pair. Per-cell values are single
    sampled-game estimates (512 passes); no aggregation, hence no
    error bars.}
  \label{fig:app-mantra}
\end{figure}


\end{document}
